\documentclass[11pt,a4paper]{ltjsarticle}
\usepackage[margin=24mm]{geometry}
\usepackage{amsmath,amssymb,amsthm,mathtools,mathrsfs}
\usepackage{booktabs,longtable,array,enumitem,needspace}
\usepackage[haranoaji]{luatexja-preset}
\usepackage[unicode,hidelinks]{hyperref}
\setlength{\emergencystretch}{2em}
\renewcommand{\proofname}{証明}
\newtheorem{theorem}{定理}[section]
\newtheorem{proposition}[theorem]{命題}
\newtheorem{lemma}[theorem]{補題}
\newtheorem{corollary}[theorem]{系}
\theoremstyle{definition}
\newtheorem{remark}[theorem]{注意}
\newtheorem{definition}[theorem]{定義}
\newcommand{\kC}{\mathbb C}
\newcommand{\cA}{\mathcal A}
\newcommand{\cE}{\mathcal E}
\newcommand{\qgr}{\operatorname{qgr}}
\newcommand{\coh}{\operatorname{coh}}
\newcommand{\Perf}{\operatorname{Perf}}
\newcommand{\End}{\operatorname{End}}
\newcommand{\HH}{\operatorname{HH}}
\newcommand{\Pic}{\operatorname{Pic}}
\newcommand{\Spec}{\operatorname{Spec}}
\newcommand{\Proj}{\operatorname{Proj}}
\newcommand{\tr}{\operatorname{tr}}
\newcommand{\rk}{\operatorname{rank}}
\title{非可換del Pezzo曲面の内在的条件とAS正則超曲面表示の限界\\\large 量子Fermat・符号型・Clifford order}
\author{}
\date{2026年11月19日／non-as-v5}
\begin{document}
\maketitle
\begin{abstract}
非可換del Pezzo曲面の候補となる内在的条件、標準AS四次元三次超曲面表示、Okawa--UedaのAS surface性を区別し、三つの条件の間に真の階層が現れることを明示例で示す。
次数三の量子Fermat点および非Calabi--Yau符号型二次元族では、離散torsion付き有限群商によってHochschild homologyに奇数次成分が現れ、Okawa--UedaのAS surfaceでない。
また、Hirzebruch曲面$\mathbb F_1$上の偶Clifford orderは次数三の反標準Hilbert関数を持ち、反標準環は次数一生成でありながら、その二次包絡の三次成分が十九次元となる。
これは標準四次元AS正則代数の三次超曲面表示を、構造対象を忘れたアーベル圏としても排除する。
次数一・二の二つのClifford orderについても、標準的な最小生成次数を持つ重み付きAS正則ambientによる超曲面表示は存在しない。一方、別の生成次数を持つAS正則ambientでは明示的な超曲面表示が存在する。
さらにトーリック商スタック$[\mathbb P^2/\mu_3]$では、反標準AZ-ampleなtorsion-free構造対象のrankが少なくとも三となることを証明し、構造層の非amplenessと特異な反標準模型を比較する。Hochschild障害とgraded ambientの障害を区別し、幾何学的・圏論的証明と有限計算を与える。
\end{abstract}
\tableofcontents

\section{はじめに}
\label{sec:main-introduction}

可換な滑らかなdel Pezzo曲面$X$では、反標準束$-K_X$の豊富性と
$K_X^2=d$が、反標準環のHilbert関数を決める。
特に$d=3$なら反標準埋込みは$\mathbb P^3$内の三次超曲面である。
非可換の場合、これに対応する\emph{内在的な曲面の条件}、
\emph{AS正則なambientを持つ超曲面表示}、
\emph{Artin--Schelter surfaceであること}は区別しなければならない。
本稿では、これらの条件の強さが異なることを、反標準自己同値を備えた
明示的な非可換del Pezzo曲面を用いて示す。

\subsection{内在的な非可換del Pezzo曲面の候補定義}

次の定義は、幾何学的な構造を表すための\emph{候補}であって、
分類定理の結論を先取りしたものではない。
$\mathcal A$は$\mathbb C$線型noetherianかつExt-finiteなアーベル圏とし、
その有限長対象によるGabriel--Krull filtrationを考える。

\begin{definition}[内在的な次数$d$の候補]\label{def:intrinsic-dp}
対$(\mathcal A,\mathcal O)$が\emph{内在的非可換del Pezzo候補（次数$d$）}
であるとは、次を満たすこととする。
\begin{enumerate}[label=(\roman*)]
\item $D^b(\mathcal A)$はSerre関手$S$を持ち、
$s:=S^{-1}[2]$が$\mathcal A$に制限される完全自己同値である。
\item $\mathcal O\in\mathcal A$はGabriel--Krull rank一かつtorsion-freeである。
すなわち、次元一以下の対象で割った商でrank一を持ち、
次元一以下の対象を非零部分対象に持たない。
\item $(\mathcal O,s)$はArtin--Zhangの意味でampleな対である
\cite{AZ1994}。
\item 反標準orbit環
\[
 R(\mathcal A,\mathcal O,s)
 :=\bigoplus_{n\geq0}\operatorname{Hom}_{\mathcal A}
       (\mathcal O,s^n\mathcal O)
\]
のHilbert級数は、可換な次数$d$のdel Pezzo曲面と同じく
\begin{equation}\label{eq:intrinsic-hilbert}
 H_R(t)=\sum_{n\geq0}\left(1+\frac d2n(n+1)\right)t^n
 =\frac{1+(d-2)t+t^2}{(1-t)^3}
\end{equation}
である。
\end{enumerate}
\end{definition}

定義\ref{def:intrinsic-dp}は、構造対象のexceptionality、
$\mathcal A$の大域次元二、導来圏のsmooth proper性をそれ自体では要求しない。
本稿の次数三Clifford例は、これらを加え、
負の次数の切断も指定した強い条件
（定義\ref{def:weak-cubic}）を満たす。
従って、以下で得る反例は内在的候補の条件が弱すぎることによる
形式的な退化ではない。
一方、Hilbert級数の定数項から
$\operatorname{End}_{\mathcal A}(\mathcal O)=\mathbb C$は既に従うが、
高次Ext群の消滅は従わない。

\subsection{三つの条件と真の包含}

次数三に限り、定義\ref{def:intrinsic-dp}を満たす対を少なくとも一つ
持つアーベル圏のクラスを$\mathfrak I_3$と書く。
そのうち、構造対象と反標準自己同値に適合する
\emph{標準AS四次元三次超曲面表示}を持つものを
$\mathfrak H^{\mathrm{hyp}}_3$とする。
具体的には、連結次数付き二次Koszul AS正則代数$G$で
大域次元四、Hilbert級数$(1-t)^{-4}$、
最小自由分解型$(1,4,6,4,1)$を持つもの、および
左右正則な三次の正規元$g\in G_3$に対し、
$R(\mathcal A,\mathcal O,s)\cong G/(g)$となる表示を要求する。
ここで圏と対はこの表示に適合しているものに限る。
この条件は、単に数値Hilbert級数を課すより強い。
二次包絡のAS正則性は、このHilbert級数と一次生成性の下では、
命題\ref{prop:quadratic-envelope-recognition}により
超曲面表示の必要十分条件となる。ただし後の$\mathcal H_3$は
さらに例外性などの圏論的条件を伴う。

さらに、Okawa--Ueda\cite[Definition~1.1]{OU}による
三次del Pezzo型のAS surfaceのクラスを
$\mathfrak{AS}^{\mathrm{OU}}_3$と書き、
\[
 \mathfrak H^{\mathrm{hyp,OU}}_3
 :=\mathfrak H^{\mathrm{hyp}}_3\cap
   \mathfrak{AS}^{\mathrm{OU}}_3
\]
と置く。OUのAS surfaceとは、
三次元AS正則$\mathbb Z$-algebraの$\operatorname{qgr}$として
得られる曲面圏であって、
本稿の「四次元AS正則ambientの三次商」とは
\emph{別の定義}である。

\begin{proposition}[二次包絡による標準三次超曲面の認識]
\label{prop:quadratic-envelope-recognition}
$k$を標数零の体とし、$R=\bigoplus_{n\ge0}R_n$を連結次数付き代数で、
$R_1$で生成され、$H_R(t)=(1-t^3)/(1-t)^4$を満たすとする。
二次包絡
\[
 B=Q(R)=T_k(R_1)/\langle\ker(R_1\otimes R_1\longrightarrow R_2)\rangle
\]
が、左右noetherianな二次Koszul AS正則代数であり、
大域次元四、Hilbert級数$(1-t)^{-4}$を持つなら、
ある三次の左右正則正規元$g\in B_3$について
\begin{equation}\label{eq:envelope-recognition}
 R\cong B/(g),\qquad \ker(B\twoheadrightarrow R)=Bg=gB
\end{equation}
となる。$B$および直線$kg\subset B_3$は$R$から関手的に決まる。
逆に、上述の型の$B$と左右正則正規元$g\in B_3$から得られる
$R=B/(g)$は、二次包絡が$B$であり、このHilbert級数を持つ。
\end{proposition}
\begin{proof}
AS正則性とHilbert級数から$\operatorname{GKdim} B=4$である。
Levasseur\cite[Theorem~6.2(ii)]{Levasseur1992}により$B$は整域となる。
自然な全射$B\twoheadrightarrow R$の核を$I$とすれば、
$B_2\cong R_2$、$\dim B_3=20$、$\dim R_3=19$なので
$I_3=kg$となる。さらに$\dim B_4=35$、$\dim R_4=31$から
$\dim I_4=4$である。整域性のため、
$gB_1$と$B_1g$はいずれも$I_4$に含まれる四次元部分空間であり、
従って$gB_1=B_1g=I_4$となる。
$B$は次数一で生成されるので$gB=Bg$であり、
整域性から$g$は左右正則である。
正規元による商のHilbert級数は
$H_{B/(g)}=(1-t^3)H_B=H_R$だから、
$B/(g)\twoheadrightarrow R$は同型となる。
最後に$B$は二次関係の核から、$kg$は$B_3\to R_3$の核から復元される。
逆向きは三次元の関係を加えても二次関係は変わらないことと、
正則元による短完全列から従う。
\end{proof}

\begin{remark}[内在条件に何を追加すべきか]
\label{rem:envelope-vs-category}
上の認識は純粋に次数付き環についての必要十分条件であり、
任意のExt-finiteアーベル圏がこの条件を自動的に満たすという主張ではない。
Artin--Zhangの再構成により、適切なnoetherian条件の下では
$\mathcal A\simeq\operatorname{qgr} R$となるが、
二次包絡のAS正則性はそれとは独立の追加条件である。
本命題は分類稿\cite[三次超曲面性の自動性]{Current}で用いた
次数三・四の議論を本稿に独立に記載したものである。
\end{remark}

\begin{proposition}[本稿の例が与える真の階層]\label{prop:three-levels}
標識を忘れたアーベル圏のクラスとして
\begin{equation}\label{eq:three-levels}
 \mathfrak H^{\mathrm{hyp,OU}}_3
 \ \subsetneq\ 
 \mathfrak H^{\mathrm{hyp}}_3
 \ \subsetneq\
 \mathfrak I_3
\end{equation}
が成り立つ。
\end{proposition}
\begin{proof}
包含自体は定義から従う。
古典的な滑らかな三次曲面は三つのクラスに入る。
一方、符号型族の三次商は、命題\ref{prop:sign-H3}により
反標準構造対象のgeneric rankが一の標準AS四次元超曲面であるが、
定理\ref{thm:sign-nonAS}により奇数次Hochschild homologyが非零で、
OUのAS surfaceではない。
従って第一の包含は真である。
量子FermatのHochschild障害は定理\ref{thm:not-OU-AS}でも独立に得られる。
次数三のClifford orderは、定理\ref{thm:main}により
内在的候補の条件（より強い条件さえ）を満たすが、
構造対象を忘れても標準graded AS四次元ambientの正規三次商ではない。
従って第二の包含も真である。
\end{proof}

\begin{remark}[包含の射程]\label{rem:ou-hyp-scope}
式\eqref{eq:three-levels}の最小のクラスは
$\mathfrak H^{\mathrm{hyp}}_3$と
$\mathfrak{AS}^{\mathrm{OU}}_3$の\emph{共通部分}である。
本稿はOUのAS surfaceがすべて標準AS四次元超曲面として
表示できると主張しない。
特に、非周期的AS正則$\mathbb Z$-algebraを許した一般のambient
認識問題は、連結graded ambientの場合とは分けて扱う。
\end{remark}

\subsection{トーリック商スタックが示す境界}\label{subsec:stack-boundary}

\begin{remark}[rank三のample構造対象と特異な反標準模型]
\label{rem:mu3-stack}
$\zeta^3=1$を原始根とし、$\mu_3$を
\[
 (x:y:z)\longmapsto(\zeta x:\zeta^{-1}y:z)
\]
によって$\mathbb P^2$に作用させる。
滑らかなトーリックDeligne--Mumfordスタック
$\mathfrak X=[\mathbb P^2/\mu_3]$について
$\mathcal A=\operatorname{coh}\mathfrak X$と置く。
このとき$D^b(\mathcal A)$のSerre関手は
$S=(-\otimes\omega_{\mathfrak X})[2]$であり、
$s=-\otimes\omega_{\mathfrak X}^{-1}$である。
$\omega_{\mathfrak X}^{-1}$は自然な線形化を持つ
$\mathcal O_{\mathbb P^2}(3)$に対応する。

スタックの構造層$\mathcal O_{\mathfrak X}$は
\emph{AZ-ampleな構造対象ではない}。
実際、三つの固定点の慣性群は$\mu_3$で、
反標準束の各冪の固定点ファイバーには自明指標しか現れない。
従って$\mathcal O_{\mathfrak X}$の負の反標準twistの直和から
非自明な指標を持つ固定点のskyscraper層への全射は存在せず、
ample pairの生成条件に反する。

これに対し、正則表現を用いた局所自由対象
\[
 \mathcal E=\mathcal O_{\mathbb P^2}\otimes_{\mathbb C}
             \mathbb C[\mu_3]
 \quad\text{（同変層として解釈する）}
\]
はgeneric rank三でtorsion-freeであり、
$(\mathcal E,s)$はAZ-ampleである。
これは$\mathcal O(3)$が粗い商に降下するample線束であり、
同変Serreの生成と消滅に加えて、
$\mathbb C[\mu_3]$が全ての既約指標を含むためである。
ただし
$\operatorname{End}_{\mathcal A}(\mathcal E)\cong\mathbb C^3$
であり、$\mathcal E$はexceptionalではない。
$\dim\operatorname{Hom}(\mathcal E,s^n\mathcal E)
=3\binom{3n+2}{2}$なので、式\eqref{eq:intrinsic-hilbert}とも異なる。

一方、$\mathcal O_{\mathfrak X}$による反標準環そのものは
\[
 \bigoplus_{n\geq0}
 H^0(\mathbb P^2,\mathcal O(3n))^{\mu_3}
 \cong
 \mathbb C[U,V,W,T]/(UVW-T^3),
 \quad (U,V,W,T)=(x^3,y^3,z^3,xyz),
\]
と明示され、Hilbert級数は
$(1-t^3)/(1-t)^4$である。
次数$3n$の不変単項式$x^ay^bz^c$では
$a-b\equiv0\pmod3$かつ$a+b+c\equiv0\pmod3$から
$a\equiv b\equiv c\pmod3$となる。共通の剰余を$r\in\{0,1,2\}$とすれば、
その単項式は$U^uV^vW^wT^r$となるので、上の表示は全次数で成立する。
この三次曲面のヤコビ行列は$(VW,UW,UV,-3T^2)$であり、
特異点は三つの座標頂点のみである。
各点の局所方程式は$vw=t^3$であるから$A_2$型となる。
従って反標準模型は特異である。
この例はrank一とHilbert級数の条件を両方外せば
rank三のample構造対象を持つスタックも候補に入ること、
しかも命題\ref{prop:mu3-minimal-rank}により、
この偏極に関してはtorsion-freeなample構造対象のrankは
三以上でなければならない。
従ってrank一・Hilbert級数・AZ-amplenessを一つずつ
吟味する必要がある。
\end{remark}

\begin{proposition}[反標準AZ-ample対象の最小rank]
\label{prop:mu3-minimal-rank}
注意\ref{rem:mu3-stack}の$\mathfrak X=[\mathbb P^2/\mu_3]$と
$s=-\otimes\omega_{\mathfrak X}^{-1}$について、
$\mathcal F\in\operatorname{coh}\mathfrak X$をtorsion-freeとする。
$(\mathcal F,s)$がArtin--Zhang ampleなら
\[
 \operatorname{rank}_{\mathfrak X}(\mathcal F)\ge3.
\]
この下界は$\mathcal E=\mathcal O_{\mathbb P^2}\otimes\mathbb C[\mu_3]$
（同変層）によって達成される。
\end{proposition}
\begin{proof}
$\mathcal F^{\vee\vee}$は滑らかな二次元Deligne--Mumfordスタック上の
反射的層であるから局所自由であり、
$Q=\mathcal F^{\vee\vee}/\mathcal F$は零次元に台を持つ。
局所環が正則二次元であるため、任意の線束$M$について
$\mathcal H\!om(Q,M)=\mathcal E\!xt^1(Q,M)=0$である。
従って任意の整数$n$で
\begin{equation}\label{eq:mu3-extension}
 \operatorname{Hom}(\mathcal F^{\vee\vee}\otimes\omega_{\mathfrak X}^{n},M)
 \xrightarrow{\ \sim\ }
 \operatorname{Hom}(\mathcal F\otimes\omega_{\mathfrak X}^{n},M)
\end{equation}
となる。これは全ての射が二重双対まで延長することを意味する。

$\mathbb P^2$の三つの座標固定点の一つ$p$を取る。
その残留慣性群は$\mu_3$であり、反標準束の$p$での指標は自明である。
実際、作用の重みは$(1,-1,0)$であり、$\mathcal O(3)$の
固定点ファイバーは各点で重み$0\pmod3$を持つ。
$\chi\in\widehat{\mu_3}$について、線束
$M_\chi=\mathcal O_{\mathbb P^2}\otimes\chi$を取る。
ample性の生成条件により、ある$n_i\gg0$について
$\bigoplus_i s^{-n_i}\mathcal F\twoheadrightarrow M_\chi$が存在する。
式\eqref{eq:mu3-extension}でその成分を$\mathcal F^{\vee\vee}$へ延長し、
$p$の残留gerbeへ制限する。右辺のファイバーは一次元表現$\chi$、
$s^{-n_i}$のファイバーは自明表現である。
従って上の全射が可能となるには、
$\mathcal F^{\vee\vee}|_p$が$\chi$を含まねばならない。
これは三つの$\chi$すべてについて成り立つので、
$\operatorname{rank}\mathcal F\ge3$である。

逆に正則表現の全ての指標を含む$\mathcal E$については、
$\mathcal O_{\mathbb P^2}(3)$のSerre生成・消滅を有限群同変化すると、
$\mathcal E(-3n)$の直和からの全射と、十分大きい$n$での
Homの全射性が従う。標数零では不変部分を取る操作が完全なので、
$(\mathcal E,s)$はAZ-ampleである。
\end{proof}

\subsection{主結果の位置と先行研究}

本稿で比較する現象には二種類ある。
量子Fermat例と非Calabi--Yau符号型二次元族は
標準四次元AS正則代数の三次商でありながら、
Hochschild homologyの奇数次成分を持つため
OUのAS surfaceではない。
これに対し$\mathbb F_1$上のClifford orderでは、
反標準環が次数一の四元で生成され、Hilbert級数が
三次超曲面と一致するにもかかわらず、
二次包絡の三次成分が$19$次元にとどまる
（標準AS四次元代数なら$20$次元）。
局所射影rank一exceptional対象の分類によって
構造対象を取り替えても非表示であることを示す。
さらに次数一・二のClifford orderでは、
可換模型の標準最小生成次数を持つ重み付きAS ambientを排除する一方、
別の最小生成次数を持つAS ambientによる超曲面表示が存在する。
この違いは、Hilbert級数だけから最小生成次数を読み取れないことを表す。

線束値二次形式からのClifford orderとconic bundleの関係は
Chan--Ingalls\cite{ChanIngalls2012}、
対応する導来圏の半直交分解は
Kuznetsov\cite{KuznetsovQuadric}による。
Hirzebruch曲面上のramification型は
Nasr\cite{Nasr2020}による既知の分類に属する。
AS正則ambientの構成には
\cite{CV2010,CKS2019,CKZ2017,LWW2014}を、
局所双対性には\cite{YZ1997}を用いる。
Hochschild homologyについては
\cite{Baranovsky,CGW,Ruan,ACH}を参照する。
本研究の貢献は、新しいorderの型の存在そのものではなく、
指定された具体例の反標準環の表示、
計算可能な二次包絡の障害、
および構造対象を忘れた非表示の証明にある。
これらの組合せの初出についてはなお網羅的な確認を要する。

以下のClifford orderの「次数」$e$は
$\dim R_n=1+en(n+1)/2$で正規化した反標準Hilbert次数である。
中心上の有理標準因子の自己交点は$K_{\mathcal A}^2=e/4$であり、
これを反標準Hilbert次数と混同しない。
また本稿の三次Clifford反例で排除するambientは
連結gradedな標準quadratic AS正則代数であって、
任意の非周期的AS正則$\mathbb Z$-algebraに
無条件に拡張するものではない。

\section{OUのAS surfaceと判定条件}
Okawa--Ueda\cite[Definition~1.1(1), \S3]{OU}のAS surfaceとは，
3次元AS-regular $\mathbb Z$-algebra $A$に対し$\qgr A$と同値なAbel圏である．
その射影代数の標準対象は非常に強いhelixをなし，有限個の連続した項が
導来圏の完全強例外列を与える\cite[(3.9)--(3.10)]{OU}．
三次曲面型（del Pezzo次数3）の場合には，その長さは$12-3=9$である．
従って導来Morita不変性と完全例外列に対するHochschild homologyの加法性から，
\begin{equation}\label{eq:ou-odd-zero}
 \HH_i(D^b(\qgr A))=0\qquad(i\ne0)
\end{equation}
がOUの意味のAS surfaceの必要条件となる．
実際，例外対象一つが生成するdg圏のHochschild homologyは$k$が次数0にあるだけである．

\section{量子Fermat環と二次の有限群grading}
以下$k=\mathbb C$とし，$\zeta$を原始三乗根とする．
\[
\begin{aligned}
 B&=k\langle t,x,y,z\rangle/\bigl(
 tx-\zeta xt,\ ty-\zeta yt,\ tz-\zeta zt,\
 xy-yx,\ xz-\zeta zx,\ yz-\zeta zy\bigr),\\
 R&=B/(t^3+x^3+y^3+z^3).
 \end{aligned}
\]
この表示は量子Fermatの認識稿\cite{Current}で使ったものと同じである．
$X\subset\mathbb P^3$を古典的な滑らかなFermat三次曲面
\[
 X=\{t^3+x^3+y^3+z^3=0\},\qquad
 S=k[t,x,y,z]/(t^3+x^3+y^3+z^3)
\]
とする．$L=(\mathbb Z/3)^2$に，四変数の重みを
\begin{equation}\label{eq:weights}
 |t|=(1,0),\quad |x|=|y|=(0,1),\quad |z|=(-1,1)
\end{equation}
で指定する．
\[
 \mu((a,b),(c,d))=\zeta^{ad}
\]
は$L$上の正規化された$2$-cocycleであり，
$S^\mu$の積を斉次元について
\[
 u_g\star v_h=\mu(g,h)u_gv_h
\]
と定める．各四変数の三乗はtwistでも変わらず，
\begin{equation}\label{eq:cocycle-identification}
 R\simeq S^\mu.
\end{equation}
実際，交換係数は
$\mu(g,h)/\mu(h,g)=\zeta^{ad-cb}$であり，
式\eqref{eq:weights}で
$q_{tx}=q_{ty}=q_{tz}=q_{xz}=q_{yz}=\zeta$，$q_{xy}=1$となる．

$H=L^\vee\simeq(\mu_3)^2$と置く．
$H$の元を$(a,b)\in(\mathbb Z/3)^2$で表せば，その$X$への作用は
\begin{equation}\label{eq:H-action}
 (a,b):(t,x,y,z)\longmapsto
 (\zeta^at,\zeta^bx,\zeta^by,\zeta^{-a+b}z).
\end{equation}
これは$X$上の忠実な射影作用である．

\begin{proposition}[離散torsionを持つ有限商stackによる表示]
\label{prop:stack-equivalence}
$H$の非退化な$2$-cocycle $\alpha$を適切に選べば，標準enhancementの導来Morita同値
\begin{equation}\label{eq:stack-eq}
 D^b(\qgr R)\simeq D^b(\coh^\alpha[X/H])
\end{equation}
が成り立つ．$\coh^\alpha[X/H]$は$\alpha$-projective $H$-equivariant coherent層の圏である．
\end{proposition}
\begin{proof}
$L$のtwisted group algebra $T=k_\mu L$を考える．その基底を$e_g$，積を
$e_ge_h=\mu(g,h)e_{g+h}$と書く．非退化な交換形式から
$T\simeq M_3(k)$である．
$H$を$S\otimes T$へ
\[
 h(s_g\otimes e_f)=\chi_g(h)\chi_f(h)^{-1}(s_g\otimes e_f)
\]
で作用させると，式\eqref{eq:cocycle-identification}から
\[
 (S\otimes T)^H\simeq S^\mu=R.
\]
$C=(S\otimes T)\# H$という通常のskew-group algebraを取り，
$e=|H|^{-1}\sum_{h\in H}h$をReynolds冪等元とする．すると
$eCe\simeq R$である．部分代数$T\# H\subset C$は
$T$の左正則表現と$H$の重み指標作用から
$\operatorname{End}_k(T)\simeq M_9(k)$に一致する．
$e$はこの$M_9(k)$の非零冪等元なのでfullであり，$C$でもfullである．
よって$C$と$R$は$\mathbb Z$次数を保つMorita同値を持つ．

一方，$H$の$T=M_3(k)$への作用はinnerであり，それを実現する
$3$次元projective表現の乗法のずれを$\alpha$とする．
この$\alpha$の交換形式は非退化である．
行列因子の内的作用を取り除けば，$C$は
$S\#_\alpha H$の行列環と同型になる（$\alpha$またはその逆への選択は規約による）．
従って$\operatorname{gr}R$と$\operatorname{gr}(S\#_\alpha H)$が同値である．
上のMorita同値は全て次数0の有限次元双加群で与えられるので，
有限次元torsionを保つ．商圏$\qgr$へ降下し，有限群equivariant Serre対応から
式\eqref{eq:stack-eq}が従う．
この構成は離散torsion付きcrossed productの標準的な枠組み
\cite{CGW}を，ここでのtwisted group algebraに具体化したものである．
\end{proof}

\section{固定楕円曲線と奇数次Hochschild homology}
有限群商のorbifold HKRによれば，捩れ付きsectorは
\begin{equation}\label{eq:twisted-hkr}
 \HH_n(D^b(\coh^\alpha[X/H]))
 \simeq \bigoplus_{g\in H}
 \left(\bigoplus_{p-q=n}H^q(X^g,\Omega^p_{X^g})
 \otimes\chi_g\right)^H,
\end{equation}
で与えられる．$\chi_g(h)=\alpha(h,g)/\alpha(g,h)$は離散torsionの推移指標である．
式\eqref{eq:twisted-hkr}は，twisted crossed productのHochschild鎖複体を\(\bigoplus_{g\in H}\)個のg-捩れ双加群成分に分解し，有限群不変部分を取って得られる．非単位元gに対しては固定点集合の法線方向で1-gが可逆なので，局所Koszul複体は固定部分に縮約し，通常のHKRが適用できる．群の共役作用での$g$成分の係数が$\chi_g$である．通常のorbifoldの公式はBaranovsky\cite{Baranovsky}，
離散torsionの指標補正は\cite{CGW,Ruan}のcrossed-productおよびinner local systemの公式に対応する．

$g=(a,b)\in H$の非単位元の固定点集合を調べる．
\begin{center}\small
\begin{tabular}{ccl}
\toprule
$g$&$X^g$&幾何\\\midrule
$(0,1),(0,2)$&$E_t=\{t=0\}\cap X$&滑らかな楕円曲線\\
$(1,1),(2,2)$&$E_z=\{z=0\}\cap X$&滑らかな楕円曲線\\
$(1,0),(2,0)$&3点&$t=z=0$のFermat零点\\
$(1,2),(2,1)$&6点&$t=z=0$と$x=y=0$の各3点\\\bottomrule
\end{tabular}
\end{center}
これは式\eqref{eq:H-action}の四座標の固有値の一致する組から従う．

\begin{theorem}[OUの意味のAS surfaceではない]
\label{thm:not-OU-AS}
量子Fermat三次曲面$\mathcal A=\qgr R$について
\begin{equation}\label{eq:HHdims}
 \boxed{\dim_k\HH_{-1}(D^b\mathcal A)=2,\qquad
 \dim_k\HH_0(D^b\mathcal A)=5,\qquad
 \dim_k\HH_1(D^b\mathcal A)=2.}
\end{equation}
従って$\mathcal A$はOU\cite[Definition~1.1]{OU}のAS surfaceではなく，
とくにOUの三次del Pezzo型AS surfaceでもない．
\end{theorem}
\begin{proof}
$E_t$では$h=(a,b)\in H$は射影的に
$z\mapsto\zeta^{-a}z$として作用する．
従って$H^0(E_t,\Omega^1)$の指標は$\zeta^{-a}$，
$H^1(E_t,\mathcal O)$の指標は$\zeta^a$である．
同様に$E_z$では$t\mapsto\zeta^{a-b}t$なので，
$H^0(E_z,\Omega^1)$の指標は$\zeta^{a-b}$，
$H^1(E_z,\mathcal O)$の指標は$\zeta^{-a+b}$である．

$\alpha$の向きを選べば式\eqref{eq:twisted-hkr}の推移指標は
\[
 \chi_{(c,d)}(a,b)=\zeta^{ad-cb}
\]
となる．$E_t$の二つのsector $g=(0,1),(0,2)$のうち，
一方では$H^0(\Omega^1)\otimes\chi_g$が自明で，
他方では$H^1(\mathcal O)\otimes\chi_g$が自明となる．
どちらも1次元である．$E_z$の二つのsectorでも同様である．
その他の固定点集合は零次元なので$\HH_{\pm1}$には寄与しない．
単位sectorの$X$は有理曲面なので$\HH_{\pm1}(X)=0$．
従って$\HH_{-1}$と$\HH_1$はいずれも2次元である．

$\HH_0$も求める．単位sectorでは
$H^0(\mathcal O_X)^H$と$H^2(\Omega_X^2)^H$が各1次元．
$H^1(\Omega_X^1)^H$は$H^2(X,k)^H$と同じ次元で，
Lefschetz固定点公式により
\[
 \frac19\sum_{g\in H}\tr(g\mid H^2(X,k))
 =\frac19\left(7+4(-2)+2(1)+2(4)\right)=1.
\]
ここでは$\tr(g\mid H^2)=\chi(X^g)-2$，$\dim H^2(X,k)=7$を使った．
従って単位sectorは3次元寄与する．
楕円sectorの$H^0(\mathcal O)$と$H^1(\Omega^1)$は
幾何的に自明な指標を持つが，$\chi_g$は非自明なので寄与しない．
$(1,0),(2,0)$の3固定点は$H$に個々に固定されるため，非自明な$\chi_g$によって寄与が零．
$(1,2),(2,1)$の6固定点では，$t=z=0$の3点は寄与せず，
$x=y=0$の3点が$H/\langle g\rangle$の推移的な軌道をなす．
$\chi_g$は$\langle g\rangle$上自明なので，この軌道から各1次元の不変部分を得る．
従って$\dim\HH_0=3+1+1=5$．

式\eqref{eq:ou-odd-zero}より，OUのAS surfaceなら$\HH_{\pm1}=0$でなければならない．
これは式\eqref{eq:HHdims}に反する．
\end{proof}

\begin{remark}[解釈と限界]
この障害はOUのAS-regular $\mathbb Z$-algebraの次元3と，
量子FermatのambientのAS-regular環の次元4とを取り違えたものではない．
計算対象は曲面の導来圏$D^b(\qgr R)$そのものである．
大域次元2，反標準Serre関手，Artin--Zhang豊富性，構造対象の例外性と，

完全強例外列の存在は別の条件である．
この例では後者は成立しない．
同じFermat曲面の\emph{全対角群}$(\mu_3)^3$による未捩れ商stackに
完全強例外列があること\cite{IU}とも矛盾しない．
本稿の捩れ付き商に用いる$H=(\mu_3)^2$はこの全対角群とは異なる：
離散torsionの非自明な指標が楕円sectorの奇数次類を生み出している．
本稿の捩れ付きequivariantモデルは通常のderived Morita同値であり，
標準のSerre作用を保つが，構造対象を固定したweak markingとの比較には
別途その像を指定する必要がある．
\end{remark}


\section{非Calabi--Yau符号型の二次元族}\label{sec:sign-family}

\subsection{正規三次商と幾何学的条件}
四変数量子多項式環
\[
 B_{\rm s}=\mathbb C\langle x_0,x_1,x_2,x_3\rangle/
 (x_ix_j-q_{ij}x_jx_i\mid i<j)
\]
で$q_{01}=1$，その他の五つの$q_{ij}$を$-1$とする．
\[
 f_3=a x_0^3+b x_0^2x_1+c x_0x_1^2+d x_1^3,\quad
 G=f_3+x_0x_2^2+x_1x_3^2,\quad R_{\rm s}=B_{\rm s}/(G)
\]
と置き，
\begin{equation}\label{eq:sign-smooth}
 ad\Delta(f_3)\ne0,\qquad
 \Delta=b^2c^2-4ac^3-4b^3d-27a^2d^2+18abcd
\end{equation}
を仮定する．これは
\[
 X_f=\{f_3(x_0,x_1)+x_0x_2^2+x_1x_3^2=0\}\subset\mathbb P^3
\]
の滑らかさと同値である．偏微分$x_0x_2=x_1x_3=0$を用いれば，
$x_0x_1\ne0$では特異点は$f_3$の重根に対応する．
$x_0=0$または$x_1=0$の残余の場合は$a,d\ne0$が排除する．

\begin{proposition}[符号型のH3認識]\label{prop:sign-H3}
式\eqref{eq:sign-smooth}の下で$R_{\rm s}$は左右noetherianな整域，
その$\qgr$は大域次元2を持ち，Serre関手は$(-1)[2]$，
$\pi R_{\rm s}$は例外対象で，$(\pi R_{\rm s},(1))$はAZ ampleである．
Hilbert級数は$(1-z^3)/(1-z)^4$．
$B_{\rm s}$は非CYで，$G$は正規だが非中心である．
これらは\cite{Current}の強いintrinsic H3対象に属する．
\end{proposition}
\begin{proof}
$B_{\rm s}$はPBW基底を持つKoszul AS-regular整域で，
Nakayama自己同型と$G$の正規化自己同型は
\[
 \mu_{B_{\rm s}}=\sigma_G=\operatorname{diag}(1,1,-1,-1)
\]
で一致する．従って$G$は正則正規，正規商公式から
$\mu_{R_{\rm s}}=1$，AS-Gorenstein parameterは1である．
以下の有限群Morita同値により，滑らかな三次曲面の有限捩れ同変層圏と
$\qgr R_{\rm s}$は同値になるので大域次元は2である．
Serre双対の標準的なAS-Gorenstein計算から$(-1)[2]$を得る．
この点は同変模型でも確認できる：
$h=(a,b,c)\in H_0$に対して四変数作用の行列式は
$(-1)^{2a+b+c}=(-1)^a$で，定義多項式$G_{\rm comm}$への
指標も$(-1)^a$である．従ってcanonical bundleの余分な
同変指標は$\det(h)/\chi_G(h)=1$となる．

整域性は省略しない．$x_1$は正則正規である：
$B_{\rm s}/(x_1)$が整域で$G\bmod x_1=ax_0^3+x_0x_2^2\ne0$
だから，$x_1$の左右消去ができる．$x_1$をOre局所化し，次数零部分で
$y=x_0x_1^{-1},u=x_2x_1^{-1},v=x_3x_1^{-1}$と書くと
\[
 C=\mathbb C[y]\langle u,v\rangle/
 (uv+vu,\ v^2+yu^2-f_3(y,1)).
\]
$s=u^2$は中心で，$C$は$\mathbb C[y,s]$上$1,u,v,uv$を基底とする．
$K=\mathbb C(y,s)$上のquaternion algebra
$(s,f_3(y,1)-ys)$を考えると，$w=f_3(y,1)-ys=0$の
付値で$s$の剰余は$f_3(y,1)/y$となる．$d\ne0$から
$y=0$に単純な極を持つので非平方である．この二次非分岐拡大の
normの付値は偶数だが$w$の付値は1．よってquaternion algebraは
斜体であり，$C$とそのskew Laurent拡大$R_{\rm s}[x_1^{-1}]$，
従って$R_{\rm s}$は整域である．generic斜体上で構造対象はrank one．
AS-Gorenstein深さと局所コホモロジー消滅から
$\operatorname{Hom}(\pi R_{\rm s},\pi R_{\rm s}(n))=(R_{\rm s})_n$，
$\operatorname{Ext}^1(\pi R_{\rm s},\pi R_{\rm s}(n))=0$，
$\operatorname{Ext}^2(\pi R_{\rm s},\pi R_{\rm s})=0$．
標準的Serre消滅と$\chi$条件からAZ amplenessが従う．
\end{proof}

\subsection{退化twisted group algebraのMorita還元}
\begin{proposition}\label{prop:sign-Morita}
$\Gamma=(\mathbb Z/2)^3$の生成元$e_0,e_1,e_2$に対し，
$x_0,x_1$の重みを$e_0$，$x_2$を$e_1$，$x_3$を$e_2$とし，
\begin{equation}\label{eq:sign-cocycle}
 \nu(u,v)=(-1)^{u_1v_0+u_2v_0+u_2v_1}
\end{equation}
とする．$H=\Gamma^\vee$，その部分群
\[
 H_0=\{(a,b,c)\in(\mathbb Z/2)^3:a+b+c=0\}
 \simeq(\mathbb Z/2)^2
\]
を置く．$H_0$は$X_f$に
\begin{equation}\label{eq:sign-H-action}
 (a,b,c)[x_0:x_1:x_2:x_3]
 =[(-1)^a x_0:(-1)^a x_1:(-1)^b x_2:(-1)^c x_3]
\end{equation}
で作用する．$H_0$の非自明なSchur multiplierを$\beta$とすると，
標準enhancementを保つ同値
\begin{equation}\label{eq:sign-quotient}
 D^b(\qgr R_{\rm s})\simeq D^b(\coh^\beta[X_f/H_0])
\end{equation}
が成立する．対応するdiscrete torsionの推移指標は
\begin{equation}\label{eq:sign-character}
 \chi_g(h)=(-1)^{a_g b_h+b_g a_h}.
\end{equation}
\end{proposition}
\begin{proof}
$S_f=\mathbb C[x_0,x_1,x_2,x_3]/(G_{\rm comm})$とする．
$G$は$\Gamma$斉次なので式\eqref{eq:sign-cocycle}による
cocycle twistは$S_f^\nu=R_{\rm s}$を与える．
$T=\mathbb C_\nu\Gamma$の交換形式のradicalは
$\langle e_0+e_1+e_2\rangle$．この次数の基底元の平方は$-1$なので
\[
 T\simeq M_2(\mathbb C)\oplus M_2(\mathbb C).
\]
特に$T$を単一の行列環とみなしてはならない．
$H$は$T$を重み指標で作用させ，二つの中心単純成分を推移的に交換し，
その安定化群は$H_0$である．

$h(s_u\otimes e_v)=u(h)v(h)^{-1}(s_u\otimes e_v)$とすれば
$(S_f\otimes T)^H=R_{\rm s}$．$T\#H$は左正則表現と指標作用により
$\operatorname{End}_{\mathbb C}(T)\simeq M_8(\mathbb C)$なので，
Reynolds冪等元は$(S_f\otimes T)\#H$でfullである．
よって$R_{\rm s}$とこのskew-group algebraは次数付きMorita同値．
二つの中心冪等元の一方のcornerに縮約すると
$(S_f\otimes M_2)\#H_0$とのMorita同値を得る．
$H_0$の$M_2$作用はinnerであり，行列因子を除去すれば
$S_f\#_\beta H_0$に帰着する．
$H_0$は$M_2$へ忠実に内的に作用するので$\beta$は
Klein四元群の唯一の非自明類で，その交換指標が
式\eqref{eq:sign-character}である．
全双加群は次数零で有限生成，torsionを保つから$\qgr$へ降下する．
有限群equivariant Serre対応によって式\eqref{eq:sign-quotient}を得る．
\end{proof}

\subsection{Hochschild障害}
\begin{theorem}[符号型二次元族はOUのAS surfaceではない]
\label{thm:sign-nonAS}
式\eqref{eq:sign-smooth}の全ての符号型曲面について
\begin{equation}\label{eq:sign-HH}
 \boxed{\dim\HH_{-1}=2,\qquad\dim\HH_0=4,\qquad\dim\HH_1=2}
\end{equation}
が成り立つ．従ってOUの意味のAS surfaceではない．
\cite{Current}の固定符号ambientのpointed分類により，
これらの一般同型類は二次元の族をなす．
\end{theorem}
\begin{proof}
$H_0=\{1,g_A,g_B,g_C\}$，
\[
 g_A=(0,1,1),\quad g_B=(1,0,1),\quad g_C=(1,1,0)
\]
とする．固定集合は
\[
\begin{array}{c|c|c}
 g& X_f^g&\chi_{\rm top}(X_f^g)\\\hline
 g_A& L_{23}=\{x_0=x_1=0\}\simeq\mathbb P^1
       \ \sqcup\ \{x_2=x_3=0,\ f_3=0\}\ (3\text{点})&5\\
 g_B& E_2=\{x_2=0\}\cap X_f\ \sqcup\{[0:0:1:0]\}&1\\
 g_C& E_3=\{x_3=0\}\cap X_f\ \sqcup\{[0:0:0:1]\}&1
\end{array}
\]
である．$E_2,E_3$は条件\eqref{eq:sign-smooth}の下で
滑らかな楕円曲線である．例えば$E_2$の方程式は
$f_3(x_0,x_1)+x_1x_3^2=0$で，
$x_3\ne0$かつ$x_1=0$にある点も滑らか，
残りの特異性は$f_3$の重根と同値である．
$E_3$も同様．

平面三次曲線のPoincar\'e residue公式により
$H^0(E_2,\Omega^1)$への$H_0$作用は
$\det(h|_{\langle x_0,x_1,x_3\rangle})/\chi_{f_3}(h)
=(-1)^{c_h-a_h}=(-1)^{b_h}$，
$H^0(E_3,\Omega^1)$への作用は$(-1)^{c_h}$である．
これらはそれぞれ式\eqref{eq:sign-character}の
$\chi_{g_B},\chi_{g_C}$と一致する．二位数ゆえ
$H^1(E_j,\mathcal O)$でも同じ指標である．
従って式\eqref{eq:twisted-hkr}を$H_0,\beta$へ適用すると，
各楕円sectorは$\HH_{1}$と$\HH_{-1}$に各1次元寄与する．
$g_A$の固定部分は有理曲線と点だけであり，
単位sectorの$X_f$は有理曲面だから他の寄与はない．

$\HH_0$については，$X_f$の$H^2$の次元が7．
Lefschetz公式$\operatorname{tr}(g|H^2)=\chi(X_f^g)-2$により
$1,g_A,g_B,g_C$のtraceは$7,3,-1,-1$．従って
\[
 \dim H^1(X_f,\Omega^1)^{H_0}=(7+3-1-1)/4=2 .
\]
単位sectorの$\HH_0$への寄与は$1+2+1=4$．
他のsectorでは$\chi_g$が非自明．$H^0(\mathcal O)$と
$H^1(\Omega^1)$への作用は，直線・楕円曲線のいずれでも自明であり，
表にある孤立点も$H_0$に個別に固定される．
従って非単位sectorの$\HH_0$不変部分は零．
これで式\eqref{eq:sign-HH}が従い，
式\eqref{eq:ou-odd-zero}によりOU非AS性が得られる．

二次元性も説明する．一般の符号型三次関係は
$f_3(x_0,x_1)+\ell_2(x_0,x_1)x_2^2+\ell_3(x_0,x_1)x_3^2$
という八次元線型系の射影空間$\mathbb P^7$の滑らかな開集合をなす．
滑らかさから$\ell_2,\ell_3$は独立で，$\operatorname{GL}_2$で
$x_0,x_1$に正規化できる．次数一の正規直線の射影スキームは
$\mathbb P\langle x_0,x_1\rangle\sqcup\{[x_2],[x_3]\}$だから，
ambientのgraded自己同型群は
$\operatorname{GL}_2\times((\mathbb G_m)^2\rtimes\mathfrak S_2)$．
共通次数スカラーは切断の射影空間で自明に作用するので
有効群次元は$6-1=5$である．一般$f_3$の三根配置は
有限自己同型群しか持たず，残る対角作用の安定化も有限である．
従って一般商の次元は$7-5=2$．詳細な固定ambientの
幾何学的群oid分類は\cite{Current}の符号型節にある．
ここでは全intrinsicモジュライの既約成分性を主張しない．
\end{proof}

\begin{remark}[二つの現象の比較]
量子Fermat例はCY ambientと中心三次を持ち，
そのsmooth central diagonal切断はpointed同型類として一点である．
符号型は非CY ambientと非中心三次を持つ正の次元の族である．
どちらにも奇数次HHが残るため，
ambientのCY性・三次の中心性の有無のみではOU型を判定できない．
\end{remark}


\section{次数三のClifford反例：主定理と適用範囲}
\label{sec:introduction}

\subsection{内在的な三次曲面構造}

可換な滑らかな三次曲面は，反標準線形系によって
$\mathbb P^3$内の三次超曲面として実現される．
その反標準Hilbert級数は$(1-t^3)/(1-t)^4$である．
非可換曲面について，このHilbert関数と反標準自己同値に関する内在的な条件から，
同様の超曲面表示が従うかを考える．
本稿の例は，以下の条件を全て満たしながら，
標準四次元AS正則代数をambientとする三次超曲面表示を持たない．

\begin{definition}\label{def:weak-cubic}
$k=\mathbb C$とする．対$(\mathcal C,\mathcal O)$が次の条件を満たすとき，
本稿では\emph{弱標識付き非可換三次曲面}と呼ぶ．
\begin{enumerate}[label=(\arabic*)]
\item $\mathcal C$は大域次元二の$k$線型noetherianアーベル圏であり，
$D^b(\mathcal C)$はsmooth proper dg enhancementとSerre関手$S$を持つ．
\item $T=S^{-1}[2]$は$\mathcal C$の完全自己同値であり，
$(\mathcal O,T)$はArtin--Zhang ample pairである．
\item 有限長対象からなるSerre部分圏を$\mathcal C_0$とし，
$\mathcal C/\mathcal C_0$で有限長となる対象の逆像を$\mathcal C_1$とする．
$\mathcal C/\mathcal C_1$における$\mathcal O$の像は非零単純であり，
$\operatorname{Hom}(M,\mathcal O)=0$が全ての$M\in\mathcal C_1$について成り立つ．
これをGabriel--Krull rank一かつtorsion-freeという．
\item $\mathcal O$はexceptionalである．すなわち
$\operatorname{End}(\mathcal O)=k$かつ
$\operatorname{Ext}^{i}(\mathcal O,\mathcal O)=0$ ($i>0$)である．
\item 全ての整数$n$について
\[
 \dim_k\operatorname{Hom}(\mathcal O,T^n\mathcal O)
 =\begin{cases}
 1+\dfrac32n(n+1)&n\ge0,\\
 0&n<0
 \end{cases}
\]
が成り立つ．
\end{enumerate}
\end{definition}

以下，加群は右加群とする．反標準環を
\begin{equation}
 R=\bigoplus_{n\ge0}\operatorname{Hom}(\mathcal O,T^n\mathcal O),
 \qquad
 ab=T^n(a)\circ b\quad(a\in R_m,\ b\in R_n)
 \label{eq:orbit-product}
\end{equation}
と定める．この規約では，Clifford orderから得られる反標準環の積は
通常のClifford積と一致する．
次数一で生成される連結次数付き代数$R$の\emph{二次包絡}を
\[
 Q(R)=T_k(R_1)\Big/\Big\langle
 \ker(R_1\otimes R_1\longrightarrow R_2)\Big\rangle
\]
と書く．
本稿で$\mathcal H_3$と呼ぶのは，定義\ref{def:weak-cubic}の条件に加えて，
$R$が次数一で生成され，$Q(R)$が左右noetherianな標準四次元二次Koszul
AS正則代数となるクラスである．ここで「標準」とは，
Gorenstein parameterが四で，自明加群の最小自由分解が型$(1,4,6,4,1)$を持つことを指す．
従って$H_{Q(R)}(t)=(1-t)^{-4}$である．

\subsection{主結果}

\begin{theorem}\label{thm:main}
$Z=\mathbb F_1$とし，負切断を$C$，ファイバー類を$f$と書く．
\[
 C^2=-1,\qquad C f=1,\qquad f^2=0,\qquad
 L=\mathcal O_Z(C+2f),\qquad
 \mathcal V=\mathcal O_Z^{\oplus2}\oplus\mathcal O_Z(f)
\]
と置く．次数がそれぞれ$f,f,C,C+f$であるCox座標$s,t,c,d$を用い，
\[
 q:\operatorname{Sym}^2\mathcal V\longrightarrow L,
 \qquad
 q=\begin{pmatrix}
 cs^2+dt&ct^2&d\\
 ct^2&ct^2+ds&0\\
 d&0&c
 \end{pmatrix}
\]
と定める．$\mathcal A=\operatorname{Cl}_0(q)$，
$\mathcal P=\operatorname{Cl}_1(q)$とすると，次が成り立つ．
\begin{enumerate}[label=(\roman*)]
\item $(\operatorname{coh}\mathcal A,\mathcal A)$は
定義\ref{def:weak-cubic}を満たし，
\[
 T=-\otimes_{\mathcal A}\mathcal P,\qquad T^2=-\otimes L,
 \qquad \operatorname{gldim}(\operatorname{coh}\mathcal A)=2
\]
である．さらに，全ての整数$n$について
$\operatorname{Ext}^1(\mathcal A,T^n\mathcal A)=0$である．
\item 反標準環$R$は次数一の四元で生成される左右noetherian domainであり，
\[
 H_R(t)=\frac{1-t^3}{(1-t)^4},\qquad \dim_k Q(R)_3=19
\]
を満たす．特にこの弱標識付き非可換三次曲面は$\mathcal H_3$に属さない．
\item $G$を連結局所有限かつ左右noetherianな二次AS正則代数で，
大域次元四，$H_G(t)=(1-t)^{-4}$を満たすものとする．
$g\in G_3$が左右正則な正規元なら，$k$線型アーベル圏同値
\[
 \operatorname{coh}\mathcal A\simeq\operatorname{qgr}(G/(g))
\]
は存在しない．この主張において，構造対象，次数移動，偏極の保存は仮定しない．
\end{enumerate}
\end{theorem}

証明は三段階からなる．まずClifford orderの局所構造と束のコホモロジーを用いて，
内在的な条件と反標準環の次数一生成を示す．
次に六本の二次関係が三次成分に張る部分空間の次元が四十五であることを計算する．
従って二次包絡の三次成分は$64-45=19$次元となり，
標準四次元AS正則代数の三次成分に必要な二十次元に達しない．
最後に局所射影rank一exceptional対象を全て分類し，
任意のアーベル圏同値のもとで候補構造対象をこの分類に帰着させる．
四つの連続する反Serre像の間の合成を比較することで，
構造対象を指定しない非表示が従う．

\subsection{先行研究との関係と次数の規約}

線束に値を持つ二次形式，conic bundle，quaternion orderの関係は既知であり，
本稿の構成はChan--Ingalls\cite[Sections~2--3, 5--6]{ChanIngalls2012}の枠組みに属する．
本例では
\[
 \det\mathcal V=\omega_Z\otimes L^2
\]
という行列式の条件が，奇Clifford bimoduleを反標準自己同値と同一視する鍵となる．
滑らかなconic bundleのClifford成分のsmooth proper性には，
Kuznetsovの半直交分解\cite[Theorem~4.2]{KuznetsovQuadric}を用いる．
Artin--Zhangの再構成定理\cite[Theorem~4.5]{AZ1994}は，
内在的なample pairを反標準環の$\operatorname{qgr}$に結び付ける．
比較対象となるAS正則三次商の飽和性と双対性には，
Yekutieli--Zhang\cite[Theorem~4.2, Corollary~4.3]{YZ1997}を用いる．

判別曲線$D=(\det q=0)$の因子類は$3C+4f$である．
$H=C+f$を$\mathbb P^2$の直線の引き戻しとすれば
\[
 D\sim 4H-C
\]
となる．従って本例のramification型は，
四次分岐曲線上の一点で$\mathbb P^2$をblow upして得られる
既知のterminal del Pezzo orderの型である．
これはNasrの分類\cite[arXiv版，Theorems~1.1(2), 2.16]{Nasr2020}に含まれる．
本稿で扱うのは，その型に属する明示的なorderについての反標準環と，
構造対象を指定しないAS三次超曲面表示の障害である．
この具体的な表示と非表示定理の既出性は未確認である．

本稿の「次数三」は，反標準Hilbert関数における次数を指す．
orderの有理標準因子の規約\cite[arXiv版，Definition~2.1]{Nasr2020}では
\begin{equation}
 K_{\mathcal A}=K_Z+\frac12D=-\frac12(C+2f),\qquad
 K_{\mathcal A}^2=\frac34
 \label{eq:canonical-order-degree}
\end{equation}
である．$\mathcal A$の中心上の階数が四であるため，
反標準Hilbert関数の次数は$4K_{\mathcal A}^2=3$となる．
この値とorderの有理交点数$K_{\mathcal A}^2$を区別する．

定理\ref{thm:main}(iii)のambientは連結次数付き代数である．
そのassociated $\mathbb Z$-algebraによる言い換えはできるが，
任意の非周期的$\mathbb Z$-algebraをambientに許す主張は本稿の範囲に含めない．


% !TEX root = ../nc-cubic-non-as.tex
\section{Clifford order と内在的な三次曲面構造}
\label{sec:construction}

基礎体を $k=\mathbb C$ とし，加群は右加群とする．
Hirzebruch 曲面 $Z=\mathbb F_1$ の負切断を $C$，
射影 $\rho:Z\to\mathbb P^1$ のファイバー類を $f$ と書く．従って
\[
 C^2=-1,\qquad C\cdot f=1,\qquad f^2=0,
 \qquad K_Z=-2C-3f.
\]
次の線束と階数3のベクトル束を固定する．
\begin{equation}
 L=\mathcal O_Z(C+2f),\qquad
 \mathcal V=\mathcal O_Z^{\oplus2}\oplus\mathcal O_Z(f).
 \label{eq:choice-LV}
\end{equation}
$\mathbb F_1$ 上の ample cone の記述から $L$ は ample であり，$L^2=3$ である．また
\begin{equation}
 \det\mathcal V=\omega_Z\otimes L^{\otimes2}
 \label{eq:determinant-V}
\end{equation}
が成り立つ．

\subsection{二次形式と判別曲線}

$Z$ の Cox 座標を $s,t,c,d$ とし，その次数を
\[
 \deg s=\deg t=f,\qquad \deg c=C,\qquad \deg d=C+f
\]
とする．不安定集合を除いた座標空間は
$\{(s,t)\ne(0,0),\ (c,d)\ne(0,0)\}$ である．
$s,t$ のいずれかと $c,d$ のいずれかをそれぞれ1に正規化すると，
$Z$ を被覆する4つの標準アフィンチャートを得る．

$L$ に値を持つ対称形式 $q:\operatorname{Sym}^2\mathcal V\to L$ を
\begin{equation}
 q=
 \begin{pmatrix}
 cs^2+dt&ct^2&d\\
 ct^2&ct^2+ds&0\\
 d&0&c
 \end{pmatrix}
 \label{eq:explicit-q}
\end{equation}
で定める．第3行・列の束の次数は $f$ なので，
例えば $q_{13}$ の次数は $L-f=C+f$，$q_{33}$ の次数は $L-2f=C$ であり，
この行列は指定した束の間の射を定める．行列式は
\begin{equation}
 \begin{split}
 \Delta={}&c^3s^2t^2-c^3t^4+c^2ds^3+c^2dt^3\\
          &+cd^2st-cd^2t^2-d^3s.
 \end{split}
 \label{eq:explicit-discriminant}
\end{equation}
判別曲線を $D=(\Delta=0)$ と書く．その因子類は
$3L-2c_1(\mathcal V)=3C+4f$ である．

\begin{lemma}
\label{lem:smooth-discriminant}
$D$ は滑らかで連結な種数3の曲線である．
$q$ の階数は全ての点で2以上であり，付随する conic bundle
\[
 p:X=\{q=0\}\subset\mathbb P_Z(\mathcal V)\longrightarrow Z
\]
は flat で，全空間 $X$ は滑らかである．
ここで $\mathbb P_Z(\mathcal V)$ は $\mathcal V$ の直線をパラメータ化する．
\end{lemma}

\begin{proof}
各アフィンチャートで式\eqref{eq:explicit-discriminant}を正規化した多項式を
$F(u,v)$ とすると，特異点イデアル
$(F,\partial_uF,\partial_vF)\subset\mathbb Q[u,v]$ の辞書式
Gr\"obner 基底は次の通りである．
\[
\begin{array}{c|cccc}
\text{正規化}&s=c=1&s=d=1&t=c=1&t=d=1\\\hline
\text{Gr\"obner 基底}&[1]&[1]&[1]&[1]
\end{array}
\]
この計算の入力と検証方法は第\ref{sec:reproducibility}節に記す．
従って $D$ は $\mathbb C$ 上で滑らかである．
$K_Z+D=C+f$ と $H^1(Z,\mathcal O_Z(C+f))=0$ より，
Serre 双対性は $H^1(Z,\mathcal O_Z(-D))=0$ を与える．
因子の完全列から $H^0(D,\mathcal O_D)=k$ であり，$D$ は連結である．
また
\[
 2g(D)-2=D(D+K_Z)=(3C+4f)(C+f)=4
\]
だから $g(D)=3$ である．

行列の階数が1以下の点では全ての $2\times2$ 小行列式が消えるため，
行列式の微分も消える．これは $D$ の滑らかさに反する．
従って各ファイバーは非零な二次式で定義される conic であり，
Hilbert 多項式の一定性から $p$ は flat である．
$D$ 外のファイバーは滑らかである．$D$ 上では階数が2であり，
退化 conic の特異点における基底方向の微分が
判別式の非零微分に対応するので，$X$ も滑らかである．
\end{proof}

\subsection{Clifford order と一般点の階数}

Clifford 関係を $xy+yx=2q(x,y)$ と規約し，
偶部分と奇部分をそれぞれ
\[
 \mathcal A=\operatorname{Cl}_0(\mathcal V,q,L),\qquad
 \mathcal P=\operatorname{Cl}_1(\mathcal V,q,L)
\]
と書く．反対称化による束としての分解は
\begin{equation}
 \begin{split}
 \mathcal A
  &=\mathcal O_Z\oplus
     (\wedge^2\mathcal V)\otimes L^{-1}
    =\mathcal O_Z\oplus\mathcal O_Z(-C-2f)
       \oplus\mathcal O_Z(-C-f)^{\oplus2},\\
 \mathcal P
  &=\mathcal V\oplus(\det\mathcal V)\otimes L^{-1}
    =\mathcal O_Z^{\oplus2}\oplus\mathcal O_Z(f)
       \oplus\mathcal O_Z(-C-f)
 \end{split}
 \label{eq:underlying-bundles}
\end{equation}
である．積は Clifford 積であり，$\mathcal P$ は
$\mathcal A$-bimodule となる．以下
\[
 \mathcal C=\operatorname{coh}\mathcal A,\qquad
 \mathcal O=\mathcal A
\]
と置く．$\mathcal C$ は noetherian な $k$-線形アーベル圏である．

\begin{proposition}
\label{prop:division-degree-three}
一般点の代数 $\mathcal A\otimes_{\mathcal O_Z}k(Z)$ は
quaternion division algebra である．従って構造対象 $\mathcal O$ は
torsion-free で Gabriel--Krull rank 1 を持つ．
\end{proposition}

\begin{proof}
$i:D\hookrightarrow Z$ とし，
$q:\mathcal V\to\mathcal V^\vee\otimes L$ の余核を $i_*\mathcal M$ と書く．
$D$ 上での corank は1なので $\mathcal M$ は線束である．
$\mathcal K=\ker(q|_D)$，$\mathcal W=\mathcal V|_D/\mathcal K$ と置くと，
$\mathcal W$ は $L|_D$ に値を持つ非退化な階数2の二次束である．
その判別線束
\[
 \eta=\det\mathcal W\otimes L^{-1}|_D
     =\mathcal M\otimes\omega_Z|_D
\]
は $\eta^{\otimes2}\simeq\mathcal O_D$ を満たす．
符号を含めて $-\det(q|_{\mathcal W})$ を用いると，
$\mathcal O_D\oplus\eta$ は退化 conic の二つの成分を区別する
非分岐二重被覆の代数となる．
この線束には完全列
\begin{equation}
 0\longrightarrow\mathcal V\otimes\omega_Z
 \xrightarrow{q}\mathcal V^\vee\otimes L\otimes\omega_Z
 \longrightarrow i_*\eta\longrightarrow0
 \label{eq:eta-resolution-degree-three}
\end{equation}
がある．左の二つの束はそれぞれ
\[
 \mathcal O_Z(-2C-3f)^{\oplus2}\oplus\mathcal O_Z(-2C-2f),
 \qquad
 \mathcal O_Z(-C-f)^{\oplus2}\oplus\mathcal O_Z(-C-2f).
\]
後者はファイバー上の次数が $-1$ なので全コホモロジーが消える．
前者の $H^1$ は Serre 双対性と
$H^1(Z,\mathcal O_Z)=H^1(Z,\mathcal O_Z(-f))=0$ により消える．
従って $H^0(D,\eta)=0$ であり，$\eta$ は非自明である．
特に二重被覆 $\widetilde D\to D$ は連結である．

$D$ の一般点の離散付値環で $q=\langle a,b,uc\rangle$ と対角化する．
ここで $a,b,c$ は単元，$u$ は一様化元である．
偶 Clifford 代数の生成元 $I=e_1e_2,J=e_1e_3$ は
\[
 I^2=-ab,\qquad J^2=-auc,\qquad IJ=-JI
\]
を満たす．その Brauer residue は $\sqrt{-ab}$ による二次拡大であり，
上の二重被覆の関数体拡大と一致する．従って residue は非自明で，
一般点の quaternion algebra は split でない．次数2の中心単純代数は
split でなければ division である．
中心上有限なorderについて，定義\ref{def:weak-cubic}の
$\mathcal C_0$は零次元に台を持つ対象からなり，
$\mathcal C_1$は台の次元が一以下の対象からなる．
実際，零次元coherent加群は有限長であり，曲線に台を持つcoherent加群を
零次元対象で割ると，有限個の曲線の一般点での有限長加群に帰着する．
従って $\mathcal C/\mathcal C_1$ は，
一般点の division algebra の有限次元右加群圏である．
その正則加群の長さは一であり，$\mathcal A$は中心上局所自由なので
低次元に台を持つ部分加群を持たない．これで結論を得る．
\end{proof}

\subsection{局所正則性と反標準自己同値}

閉点の完備局所環を $B=k[[u,v]]$ とする．$D$ 外では
$\widehat{\mathcal A}\simeq M_2(B)$ である．$D=(u)$ 上では
単元の平方根を取って Clifford の偶生成元を正規化すると
$I^2=1,J^2=u,IJ=-JI$ とでき，
\begin{equation}
 \widehat{\mathcal A}\simeq
 \begin{pmatrix}B&B\\uB&B\end{pmatrix}
 \label{eq:local-order-degree-three}
\end{equation}
を得る．これは $u=r^2$ と置いた
$k[[r,v]]\rtimes\mu_2$ と同型であり，$\mu_2$ は $r$ の符号を変える．
群の位数が可逆なので誘導加群の平均化による分裂を用いると，
この skew group algebra の大域次元は2である．
実際，$k[[r,v]]$ 上の長さ2の射影分解を誘導して上界を得，
閉点加群の射影次元から下界を得る．
忠実平坦な完備化から射影性が降下するので，元の局所 order でも
coherent module の射影次元は2以下である．

各点の近傍で $q$ は単元を表すので，可逆な Clifford の奇元を取ることができる．
従って Clifford 積は bimodule の同型
\begin{equation}
 \mathcal P\otimes_{\mathcal A}\mathcal P\simeq\mathcal A\otimes L
 \label{eq:odd-square}
\end{equation}
を与える．また最高外冪成分への完全 pairing
\[
 \mathcal A\otimes\mathcal P\longrightarrow
       \det\mathcal V\otimes L^{-1}
\]
は偶元と奇元に関して cyclic なので，
$\mathcal A^\vee\simeq
\mathcal P\otimes(\det\mathcal V)^{-1}\otimes L$ は bimodule の同型である．
式\eqref{eq:determinant-V}より
\[
 \omega_{\mathcal A}:=
 \mathcal Hom_{\mathcal O_Z}(\mathcal A,\omega_Z)
 \simeq\mathcal P\otimes L^{-1}\simeq\mathcal P^{-1}.
\]
局所射影分解と $Z$ 上の双対性から，Serre functor は
\begin{equation}
 S=(-\otimes_{\mathcal A}\mathcal P^{-1})[2],\qquad
 T=S^{-1}[2]=-\otimes_{\mathcal A}\mathcal P,
 \qquad T^2=-\otimes L
 \label{eq:Serre-degree-three}
\end{equation}
である．従って $T$ はアーベル圏の完全自己同値である．
Serre 双対性と標準$t$-構造により，$i>2$なら
$\operatorname{Ext}^i_{\mathcal C}(M,N)=0$である．一方
\[
 \operatorname{Ext}^2_{\mathcal C}(\mathcal A,\mathcal P^{-1})
 \simeq\operatorname{End}_{\mathcal C}(\mathcal P^{-1})^*\ne0
\]
だから，$\mathcal C$のhomological dimensionは正確に二である．
Hom と Ext の有限次元性は中心 $Z$ の射影性から従う．
さらに conic bundle の半直交分解\cite[Theorem~4.2]{KuznetsovQuadric}
\begin{equation}
 D^b(X)=\left\langle
 D^b(Z,\mathcal A),\ p^*D^b(Z)\otimes\mathcal O_{X/Z}(1)
 \right\rangle
 \label{eq:conic-sod-degree-three}
\end{equation}
により，$D^b(\mathcal C)$ の標準 dg enhancement は smooth proper である．

\begin{proposition}
\label{prop:ample-degree-three}
$(\mathcal O,T)$ は Artin--Zhang ample である．反標準環
\[
 R=\bigoplus_{n\ge0}\operatorname{Hom}_{\mathcal C}
                  (\mathcal O,T^n\mathcal O)
\]
について $\mathcal C\simeq\operatorname{qgr}R$ が成り立つ．
\end{proposition}

\begin{proof}
任意の coherent $\mathcal A$-module $M$ は中心上でも coherent である．
$m\gg0$ では $M\otimes L^m$ は大域生成され，随伴により
$\mathcal A^{\oplus N}\twoheadrightarrow M\otimes L^m$ を得る．
これは $(T^{-2m}\mathcal O)^{\oplus N}\twoheadrightarrow M$ を与える．
次に全射 $M\twoheadrightarrow N$ の核を $K$ とし，$n=2m+r$，$r=0,1$ と書く．
$\mathcal A$ の局所射影性と Serre 消滅により
\[
 \operatorname{Ext}^1_{\mathcal A}(\mathcal A,T^nK)
 =H^1\bigl(Z,(K\otimes_{\mathcal A}\mathcal P^r)\otimes L^m\bigr)=0
 \qquad(m\gg0).
\]
これが ample pair のもう一つの条件である．
再構成は Artin--Zhang の定理\cite[Theorem~4.5]{AZ1994}による．
\end{proof}

\subsection{全次数の Hilbert 関数とコホモロジー消滅}

\begin{proposition}
\label{prop:Hilbert-degree-three}
構造対象 $\mathcal O$ は exceptional であり，全ての整数 $n$ に対して
\[
 \operatorname{Ext}^1_{\mathcal C}(\mathcal O,T^n\mathcal O)=0
\]
が成り立つ．$n<0$ では
$\operatorname{Hom}_{\mathcal C}(\mathcal O,T^n\mathcal O)=0$，
$n\ge0$ では
\begin{equation}
 \dim_k R_n=1+\frac32n(n+1),\qquad
 \operatorname{Ext}^{i}_{\mathcal C}(\mathcal O,T^n\mathcal O)=0
 \quad(i>0).
 \label{eq:Hilbert-degree-three}
\end{equation}
特に
\[
 H_R(t)=\frac{1-t^3}{(1-t)^4}.
\]
\end{proposition}

\begin{proof}
射影 $\rho$ に沿った計算により，$a\ge0$ に対して
\[
 \rho_*\mathcal O_Z(aC+bf)
       =\bigoplus_{i=0}^a\mathcal O_{\mathbb P^1}(b-i),
 \qquad R^1\rho_*\mathcal O_Z(aC+bf)=0.
\]
従って $b\ge a$ なら
\[
 h^0\mathcal O_Z(aC+bf)
   =\sum_{i=0}^a(b-i+1)
   =(a+1)(b+1)-\frac{a(a+1)}2,
\]
かつ高次コホモロジーは消える．
$a=-1$ の線束はファイバー上で acyclic なので全コホモロジーが消える．

式\eqref{eq:odd-square}と\eqref{eq:underlying-bundles}から
\[
 R_{2m}=H^0(Z,\mathcal A\otimes L^m),\qquad
 R_{2m+1}=H^0(Z,\mathcal P\otimes L^m).
\]
上の線束の公式を各直和因子に適用すると，$m\ge0$ の全範囲で
\begin{equation}
 \dim R_{2m}=6m^2+3m+1,\qquad
 \dim R_{2m+1}=6m^2+9m+4
 \label{eq:Hilbert-parities}
\end{equation}
を得る．$m=0$ で現れる負の線束はファイバー次数 $-1$ であり，
この場合も公式は成立する．同時に高次コホモロジーの消滅を得る．
$m<0$ では偶数・奇数のいずれの直和因子もファイバー次数が負なので，
全ての $n<0$ に対して $\operatorname{Hom}(\mathcal O,T^n\mathcal O)=0$ となる．
Serre 双対性は
\[
 \operatorname{Ext}^1(\mathcal O,T^n\mathcal O)^*
 \simeq\operatorname{Ext}^1(\mathcal O,T^{-n-1}\mathcal O)
\]
を与えるため，負の $n$ に対する $\operatorname{Ext}^1$ の消滅も従う．
$n=0$ の場合が構造対象の exceptionality である．
\end{proof}

\subsection{Hochschild homology}

\begin{proposition}
\label{prop:HH-degree-three}
標準的な dg enhancement について，HKR の次数を $p-q=i$ と規約すると
\[
 \operatorname{HH}_{-1}(\mathcal C)=k^2,\qquad
 \operatorname{HH}_0(\mathcal C)=k^4,\qquad
 \operatorname{HH}_1(\mathcal C)=k^2,
 \qquad \operatorname{HH}_i(\mathcal C)=0\quad(|i|\ge2)
\]
である．特に $D^b(\mathcal C)$ は full exceptional collection を持たない．
\end{proposition}

\begin{proof}
滑らかな判別曲線を持つ conic bundle に対する公式
\cite[Theorem~8.12]{KuznetsovHH}を上の次数規約に直すと
\[
 \operatorname{HH}_i(\mathcal C)
 \simeq\operatorname{HH}_i(Z)
 \oplus H^{-i}(D,\eta)
 \oplus H^{1-i}(D,\omega_D\otimes\eta)
\]
となる．$Z=\mathbb F_1$ の Hodge cohomology は対角部分だけで，
$\operatorname{HH}_0(Z)=k^4$ である．
$\eta$ は種数3の曲線上の非自明な2-torsion 線束なので，
Riemann--Roch と Serre 双対性は
\[
 h^1(D,\eta)=h^0(D,\omega_D\otimes\eta)=2,
 \qquad h^0(D,\eta)=h^1(D,\omega_D\otimes\eta)=0
\]
を与える．これで主張を得る．
full exceptional collection が存在すれば Hochschild homology の加法性により
非零次数の部分が消えるため，最後の結論も従う．
\end{proof}

% !TEX root = ../nc-cubic-non-as.tex
\section{反標準環と二次包絡}
\label{sec:quadratic-obstruction}

\subsection{次数一生成と環論的性質}
\label{subsec:generation-all-degrees}

$\mathcal V=\mathcal O_Z^{\oplus2}\oplus\mathcal O_Z(f)$の大域切断
\[
 x=(1,0,0),\qquad y=(0,1,0),\qquad
 z=(0,0,s),\qquad w=(0,0,t)
\]
を，$\mathcal V\hookrightarrow\mathcal P$を通して$R_1$の元とみなす．
式\eqref{eq:underlying-bundles}より，これらは$R_1$の基底である．

\begin{proposition}
\label{prop:generation-all-degrees}
$R$ は次数一の四元 $x,y,z,w$ で生成される．さらに $R$ は
有限生成可換中心部分環上の有限加群であり，左右 noetherian な domain である．
\end{proposition}

\begin{proof}
中心の切断環を
\[
 S=\bigoplus_{m\ge0}H^0(Z,L^m),\qquad
 \deg H^0(Z,L^m)=2m
\]
と置く．$\mathcal A$ のスカラー部分と Clifford 積により，
$S$ は $R$ の中心部分環である．Cox 座標での展開は
\begin{equation}
 H^0(Z,L^m)
 =\bigoplus_{j=0}^m c^{m-j}d^j\,k[s,t]_{2m-j}
 \qquad(m\ge0)
 \label{eq:scroll-monomials}
\end{equation}
である．二変数斉次多項式の積が各次数で全射になることから，
$S$ は5つの切断
\[
 cs^2,\quad cst,\quad ct^2,\quad ds,\quad dt
\]
で生成される．一方，式\eqref{eq:explicit-q}の Clifford 関係は
\begin{equation}
 cs^2=z^2,\qquad cst=zw=wz,\qquad ct^2=w^2,
 \qquad ds=\frac{xz+zx}{2},\qquad dt=\frac{xw+wx}{2}
 \label{eq:scalar-generators}
\end{equation}
を与える．従って $S$ 全体が $x,y,z,w$ の生成する部分代数に含まれる．

同様に
\[
 H^0(Z,L^m\otimes\mathcal O_Z(f))
 =\bigoplus_{j=0}^m c^{m-j}d^j\,k[s,t]_{2m+1-j}
\]
なので，次数付き $S$-加群
$\bigoplus_{m\ge0}H^0(Z,L^m\otimes\mathcal O_Z(f))$ は
最初の成分 $H^0(Z,\mathcal O_Z(f))=\langle s,t\rangle$ で生成される．
従って式\eqref{eq:underlying-bundles}の外積分解を用いると，
偶数次数部分は $S$ 上，スカラー $1$ と
\[
 x\wedge y,\quad x\wedge z,\quad x\wedge w,\quad
 y\wedge z,\quad y\wedge w
\]
で生成される．ここに現れる外積は全て次数2の元であり，
例えば $x\wedge y=(xy-yx)/2$ なので
$x,y,z,w$ の積の線形結合である．
奇数次数部分は $S$ 上，$x,y,z,w$ と
\[
 x\wedge y\wedge z,\qquad x\wedge y\wedge w
\]
で生成される．最後の二つは次数3に属し，それぞれ
$x,y,z$ と $x,y,w$ の積を全反対称化した元である．
以上より $R$ の全ての次数が $x,y,z,w$ で生成される．
同時に上の有限生成系は $R$ が $S$ 上有限であることを示す．
$S$ は有限生成可換 $k$-代数なので，$R$ は左右 noetherian である．

最後に一般点で $D_0=\mathcal A\otimes k(Z)$ と置く．
命題\ref{prop:division-degree-three}により $D_0$ は division algebra である．
可逆 bimodule $\mathcal P\otimes k(Z)$ は右 $D_0$-加群として階数1であり，
その tensor 冪の直和は skew polynomial ring と同型になる．
各 $\mathcal P^{\otimes_{\mathcal A}n}$ は中心上のベクトル束なので，
その大域切断は一般点の切断に単射する．従って $R$ はこの
skew polynomial ring の部分代数であり，domain である．
\end{proof}

\subsection{六本の二次関係と三次成分の計算}

命題\ref{prop:generation-all-degrees}より積$R_1\otimes R_1\to R_2$は全射である．
命題\ref{prop:Hilbert-degree-three}から$\dim R_1=4$，$\dim R_2=10$なので，
その核は六次元である．

\begin{lemma}\label{lem:six-quadrics}
上の核は次の独立な六本の二次関係で張られる．
\begin{align}
 r_1&=zw-wz,&
 r_2&=yz+zy,&
 r_3&=yw+wy,\notag\\
 r_4&=2x^2-2z^2-xw-wx,&
 r_5&=xy+yx-2w^2,&
 r_6&=2y^2-2w^2-xz-zx.
 \label{eq:six-quadrics}
\end{align}
$V=R_1$，$W=\langle r_1,\ldots,r_6\rangle$と置くと，二次包絡
$Q=T_k(V)/(W)$は
\[
 \dim_k Q_3=19
\]
を満たす．
\end{lemma}
\begin{proof}
六本の式がClifford関係から従うことは，
式\eqref{eq:explicit-q}と\eqref{eq:scalar-generators}から分かる．
それらの独立性は，順に$zw,yz,yw,x^2,xy,y^2$の係数を見ると従う．
核の次元が六であるからこれで核全体となる．

三次成分は
\[
 Q_3=V^{\otimes3}/(W\otimes V+V\otimes W)
\]
である．$(x_1,x_2,x_3,x_4)=(x,y,z,w)$とし，
三次単項式をこの順の辞書式順序で並べる．
$(r_jx_i)_{j,i}$，$(x_ir_j)_{j,i}$をこの順に並べた
$64\times48$整数行列を$M$とする．
$M$には行列式$512$の$45\times45$小行列が存在し，
また三本の独立な整数右核ベクトルが存在する．
小行列の添字，核ベクトル，および有理数体上の検証コードを
第\ref{sec:reproducibility}節に指定する．
したがって$\operatorname{rank}M\ge45$かつ
$\operatorname{rank}M\le48-3=45$であり，
$\dim Q_3=64-45=19$となる．
\end{proof}

命題\ref{prop:generation-all-degrees}より$Q_3\to R_3$は全射であり，
両辺の次元が十九なので同型である．すなわち，二次関係の三次帰結だけで
反標準環の三次成分の次元に達している．

\subsection{構造対象と偏極を保持した非表示}

標準的なnc $\mathbb P^3$の斉次座標環は，次数一で生成される
二次AS正則代数$G$で$H_G(t)=(1-t)^{-4}$を満たすものとする．
$g\in G_3$が正則正規元なら$B=G/(g)$について
$B_1=G_1$，$B_2=G_2$であり，$G$は二次であるから
\[
 G=T_k(B_1)/\bigl(\ker(B_1\otimes B_1\longrightarrow B_2)\bigr).
\]
従ってその二次包絡の三次成分は必ず$\dim G_3=20$である．
補題\ref{lem:six-quadrics}は，$R$がこのような三次商に次数付き同型でないことを示す．
以下では構造対象を保存しないアーベル圏同値も排除する．

% !TEX root = ../nc-cubic-non-as.tex
\section{局所射影rank一exceptional対象}
\label{sec:exceptional-classification}

\subsection{局所型と行列式}

\begin{lemma}\label{lem:regular-rank-one}
$F\in\operatorname{coh}\mathcal A$が局所射影で，一般点のquaternion division algebra上のrankが一なら，
$F$は局所的に正則右加群$\mathcal A$と同型である．
また，ある$N\in\operatorname{Pic}(Z)$に対し
\[
 \det F\simeq\det\mathcal A\otimes N^{\otimes2}
\]
となる．
\end{lemma}
\begin{proof}
判別曲線の一般点の離散付値環上では，orderは
$\mathcal A=R[I]\oplus R[I]J$，$J^2$が一様化元の単元倍という形であり，
$R[I]$は非自明な剰余二次拡大を持つ．
Jacobson radicalは$J\mathcal A$，その商は体であるから$\mathcal A$は局所環である．
有限射影加群は自由であり，一般点のrankから自由階数は一となる．

閉分岐点の中心の完備局所環を$R=k[[u,v]]$とすると，orderは
$\left(\begin{smallmatrix}R&R\\uR&R\end{smallmatrix}\right)$である．
二つの原始冪等元に対応する直既約射影を$P_1,P_2$と書けば，
完備化した$F$は$P_1^{\oplus a}\oplus P_2^{\oplus b}$と分解する．
形式的分岐因子の一般点への局所化は，分岐曲線の一般点で得た
自由加群の基底変換である．この局所化の後も$P_1,P_2$は
互いに非同型な直既約射影である．正則加群は$P_1\oplus P_2$なので，
Krull--Schmidtの一意性から$a=b=1$となる．
分岐外では完備orderは$M_2(R)$であり，中心上rank四の射影加群は正則加群と同型である．
完備化上の生成元を中心の極大イデアルを法として持ち上げ，Nakayamaの補題を用いると，
元の局所環でも$\mathcal A\to F$が全射になる．
同じrankのベクトル束の全射なので同型である．

従って遷移写像は$\mathcal A^\times$の元による左乗法である．
quaternion algebraにおける左乗法の行列式はreduced normの二乗だから，
reduced normの遷移関数で定まる線束$N$が所望の等式を与える．
\end{proof}

\begin{lemma}\label{lem:numerical-uniqueness}
局所射影division rank一の$F$について
$[F]=[\mathcal A]$ in $K_{\mathrm{num}}(\mathcal A)$なら，$F\simeq\mathcal A$である．
\end{lemma}
\begin{proof}
忘却関手の左随伴は$E\mapsto E\otimes\mathcal A$なので
\[
 \chi_Z(E,\operatorname{Forg}F)=\chi_{\mathcal A}(E\otimes\mathcal A,F)
\]
である．従って忘却関手は数値的群上に降り，$c_1(F)=c_1(\mathcal A)$を得る．
$\operatorname{Pic}(\mathbb F_1)$はtorsion-freeで$\operatorname{Pic}^0=0$だから，
$\det F\simeq\det\mathcal A$である．一方
\[
 c_1(\mathcal A)=-3C-4f,\qquad c_1(\omega_{\mathcal A})=-5C-8f,
 \qquad\det\omega_{\mathcal A}\otimes(\det F)^{-1}=\mathcal O_Z(-2L).
\]
非零な$F\to\omega_{\mathcal A}$は，division rank一により一般点で同型である．
従ってその行列式は$\mathcal O_Z(-2L)$の非零切断となるが，$L$のample性に反する．
Serre双対性から$\operatorname{Ext}^2_{\mathcal A}(\mathcal A,F)=0$を得る．
$\chi(\mathcal A,F)=1$なので$\operatorname{Hom}(\mathcal A,F)\ne0$である．
その非零写像の行列式は非零定数となるから，写像は全体で同型である．
\end{proof}

\subsection{数値的群とexceptional対象}

\begin{lemma}\label{lem:numerical-rank-four}
$K_{\mathrm{num}}(\mathcal A)_{\mathbb Q}$の次元は四であり，
$[\mathcal A],[\mathcal P],[\mathcal A(C)],[\mathcal A(f)]$はその基底である．
この順序のEuler行列は
\[
 \mathbf E=\begin{pmatrix}
 1&4&0&2\\
 1&1&-2&0\\
 -2&-1&1&-5\\
 0&1&-5&1
 \end{pmatrix},\qquad\det\mathbf E=100.
\]
\end{lemma}
\begin{proof}
前節のconic bundleを$p:X\to Z$とする．
任意の素因子$E\subset Z$について$p^{-1}E$は素因子である．
$E$が判別曲線と異なる場合は一般ファイバーが滑らかなconicであることから従う．
判別曲線の場合は，その二つの幾何学的成分が連結な非分岐二重被覆によって交換されるため，
基礎関数体上では既約となる．
一般conicへの制限
\[
 \operatorname{Pic}(X)_{\mathbb Q}\longrightarrow
 \operatorname{Pic}(X_{k(Z)})_{\mathbb Q}
\]
の核は垂直因子類で生成される．実際，一般ファイバーへの制限が自明な因子は，
有理関数の因子を差し引くことで垂直因子にできる．
$p$は相対次元一のflat射だから，垂直素因子の像は曲線である．
従って上の既約性から，核は$p^*\operatorname{Pic}(Z)_{\mathbb Q}$に等しい．
一般conicのPicard群はrank一で，相対$\mathcal O(1)$が非零類を与えるから$\rho(X)=\rho(Z)+1=3$である．
$p_*\mathcal O_X=\mathcal O_Z$かつ$R^{i>0}p_*\mathcal O_X=0$から
$H^{i>0}(X,\mathcal O_X)=0$を得る．
Lefschetzの$(1,1)$定理とHard Lefschetzにより$H^2(X),H^4(X)$はともに代数的であり，
Riemann--RochとPoincar\'e双対性から$\operatorname{rank}K_{\mathrm{num}}(X)=1+3+3+1=8$である．
conic bundleの半直交分解\cite[Theorem~4.2]{KuznetsovQuadric}により
\[
 D^b(X)=\langle D^b(Z,\mathcal A),\ p^*D^b(Z)\otimes\mathcal O_{X/Z}(1)\rangle
\]
となるので，$\operatorname{rank}K_{\mathrm{num}}(\mathcal A)=8-4=4$である．
ここでadmissibleな部分圏の包含と両随伴は数値的同値を保つため，
半直交分解は数値的Grothendieck群の直和分解を誘導する．

行列の計算は，$E=aC+bf$について
\[
 \chi(\mathcal A(E))=1+2(-a^2+2ab)+a+b,
 \qquad
 \chi(\mathcal P(E))=4+2(-a^2+2ab)+3(a+b)
\]
および$\mathcal P^{-1}=\mathcal P(-L)$を使う．
これらは線束分解と曲面Riemann--Rochから従う．
行列式が非零であることから，四つの類が基底である．
\end{proof}

\begin{proposition}\label{prop:exceptional-rank-one}
$\operatorname{coh}\mathcal A$の局所射影division rank一のexceptional対象は，ちょうど
\[
 \mathcal A\otimes M,\qquad\mathcal P\otimes M,
 \qquad M\in\operatorname{Pic}(Z)
\]
である．
\end{proposition}
\begin{proof}
補題\ref{lem:regular-rank-one}により，ある整数$\alpha,\beta$について
\[
 c_1(F)=c_1(\mathcal A)+2(\alpha C+\beta f)
\]
と書ける．前の基底について，rankと第一Chern類の条件は，ある$\lambda\in\mathbb Q$を用いて
\[
 [F]=\left(1-\frac\alpha2-\frac\beta2+\frac\lambda2,
 \lambda,\frac{\alpha-\lambda}{2},\frac{\beta-2\lambda}{2}\right)
\]
と表せる．Euler行列から
\[
 \chi(F,F)=1+\alpha^2-2\alpha\beta-2\alpha+5\lambda,
 \quad
 \chi(\mathcal A,F)=1-\frac\alpha2+\frac\beta2+\frac{5\lambda}{2}
\]
を得る．exceptionalityを用いて$\lambda$を消去すると
\[
 \chi(\mathcal A,F)=1+\frac{-\alpha^2+2\alpha\beta+\alpha+\beta}{2}
 =1+\frac{N^2+N\cdot L}{2},\qquad N=\alpha C+\beta f.
\]
$\alpha^2-\alpha$が偶数であることと左辺の整数性から，$\beta$は偶数である．
$\alpha$も偶数なら$\mathcal A(N/2)$，奇数なら$\mathcal P((N-L)/2)$が
$F$と同じrank，第一Chern類，自己Euler数を持ち，従って同じ数値的類を持つ．
この比較対象を$\mathcal A$に送る自己同値で$F$を正規化し，
補題\ref{lem:numerical-uniqueness}を適用する．
逆に表示した対象はいずれも，$\mathcal A$に自己同値を施して得られるため，
局所射影division rank一でexceptionalである．
\end{proof}

\section{構造対象を指定しない超曲面表示の排除}
\label{sec:unpointed-obstruction}

\subsection{AS三次商の標準対象}

\begin{lemma}\label{lem:AS-cubic-cohomology}
$G,g$を定理\ref{thm:main}(iii)のように取り，
$B=G/(g)$，$\mathcal B=\pi_BB$と置く．
\begin{enumerate}[label=(\roman*)]
\item 全ての整数$i,j$について
\[
 \operatorname{Hom}(\mathcal B(i),\mathcal B(j))=B_{j-i},\qquad
 \operatorname{Ext}^1(\mathcal B(i),\mathcal B(j))=0
\]
であり，$\mathcal B$はexceptionalである．
\item 任意の$M\in\operatorname{qgr}B$について，
$n\gg0$なら$\operatorname{Ext}^1(\mathcal B(-n),M)=0$である．
\item $D^b(\operatorname{qgr}B)$がSerre関手$S_B$を持つなら，
\[
 S_B(\mathcal B(i))\simeq\mathcal B(i-1)[2]
 \qquad(i\in\mathbb Z)
\]
が対象として成り立つ．従って$T_B=S_B^{-1}[2]$について
$T_B^n\mathcal B\simeq\mathcal B(n)$である．
\end{enumerate}
\end{lemma}

\begin{proof}
$G$のAS-Gorenstein parameterを$\ell$とする．
次数一生成，二次性，AS双対性から，自明加群の最小自由分解の
交代Hilbert多項式は
\[
 1-4t+6t^2-4t^{\ell-1}+t^\ell
\]
となる．これが$H_G(t)^{-1}=(1-t)^4$に等しいので$\ell=4$である．
二項自由分解$0\to G(-3)\to G\to B\to0$と
change-of-rings spectral sequenceにより，
$\underline{\operatorname{Ext}}_G^q(B,G)$は$q=1$に集中し，
その項はbimoduleの捩れを除いて$B(3)$である．
従って，左右のAS-Gorenstein条件と有限injective dimensionは商へ降り，
$B$のparameterは一，injective dimensionは三となる．

Yekutieli--Zhang\cite[Theorem~4.2, Corollary~4.3]{YZ1997}により，
$B$のbalanced dualizing complexはある次数付き自己同型$\sigma$を用いて
${}_\sigma B(-1)[3]$と書ける．局所双対性から
$H^j_{\mathfrak m}(B)=0$ ($j\ne3$)である．
局所コホモロジーと切断の完全列，およびその高次版により，
(i)のHomと$\operatorname{Ext}^1$の公式を得る．
さらに$\operatorname{cd}(\operatorname{qgr}B)=2$であり，
\[
 \operatorname{Ext}^2(\mathcal B,\mathcal B)^*
 =H^3_{\mathfrak m}(B)_0^*
 =B(-1)_0=0
\]
だから，$\mathcal B$はexceptionalである．

有限生成次数付き加群$\widetilde M$で$\pi_B\widetilde M=M$となるものを取る．
同じ局所コホモロジー比較と\cite[Theorem~4.2(1)]{YZ1997}から
\[
 \operatorname{Ext}^1(\mathcal B(-n),M)
 =H^2_{\mathfrak m}(\widetilde M)_n=0\qquad(n\gg0)
\]
となる．最後の消滅は局所コホモロジーの右有界性であり，これが(ii)である．

$\operatorname{qgr}B$上の双対化対象は$\pi_B({}_\sigma B(-1))[2]$である．
正則加群の自己同型による捩れは右加群として正則加群と同型なので，
これは$\mathcal B(-1)[2]$と同型である．
Serre関手が存在すれば，双対性を表現する対象の一意性から(iii)を得る．
この主張は対象ごとの同型であり，$T_B$と次数移動の関手としての同一視を含まない．
\end{proof}

\subsection{アーベル圏同値の排除}

\begin{theorem}\label{thm:unpointed-graded}
$G$を連結局所有限かつ左右noetherianな二次AS正則代数で，大域次元四，
$H_G(t)=(1-t)^{-4}$を満たすものとする．
$g\in G_3$を正則正規元とすると，$k$線型アーベル圏同値
\[
 \operatorname{qgr}(G/(g))\simeq\operatorname{coh}\mathcal A
\]
は存在しない．
\end{theorem}
\begin{proof}
$B=G/(g)$，$\mathcal B=\pi_B B$と置き，そのような同値$\Phi$があると仮定する．
この同値は有界導来圏の同値を誘導する．従って
$D^b(\operatorname{qgr}B)$はSerre関手$S_B$を持ち，
$\Phi$は$T_B=S_B^{-1}[2]$と$T$を可換にする．
補題\ref{lem:AS-cubic-cohomology}により
$T_B^n\mathcal B\simeq\mathcal B(n)$である．

$F=\Phi\mathcal B$とする．閉点に台を持つ任意の単純$\mathcal A$加群$Q_0$について
$T^{2m}Q_0\simeq Q_0$だから
\[
 \operatorname{Ext}^1(F,Q_0)
 \simeq\operatorname{Ext}^1(T^{-2m}F,Q_0)
 \simeq\operatorname{Ext}^1(\mathcal B(-2m),\Phi^{-1}Q_0)=0
 \qquad(m\gg0).
\]
最後は補題\ref{lem:AS-cubic-cohomology}(ii)による．
値を取る加群が零次元なのでlocal-to-global spectral sequenceは，
局所$\operatorname{Ext}^1$も零であることを示す．
完備局所orderはsemiperfectであり，最小射影表示の核が非零ならその単純商への
$\operatorname{Ext}^1$が非零となるから，$F$は完備化上で射影である．
忠実平坦降下により$F$は局所射影である．

一般点でのdivision rankを$r$とすると，
$\mathcal H om_{\mathcal A}(F,F)$の中心上のrankは$4r^2$である．
$L^2=3$，曲面Riemann--Roch，Serre消滅から
\[
 \dim\operatorname{Hom}(F,T^{2m}F)=6r^2m^2+O(m).
\]
一方，左辺は$\dim B_{2m}=6m^2+3m+1$だから$r=1$である．
$F$はexceptionalなので，命題\ref{prop:exceptional-rank-one}により
ある$M\in\operatorname{Pic}(Z)$について
$F\simeq\mathcal A\otimes M$または$\mathcal P\otimes M$である．
いずれも$T$と可換な自己同値によって$\mathcal A$から得られる．

最後に四つの対象$F,TF,T^2F,T^3F$の間の合成を考える．
連続するHom空間を$V_i=\operatorname{Hom}(T^iF,T^{i+1}F)$とし，
二段の合成の核を
\[
 W_i=\ker\bigl(V_i\otimes V_{i+1}\longrightarrow
 \operatorname{Hom}(T^iF,T^{i+2}F)\bigr)
\]
とする．二次合成関係から定まる三次空間は
\begin{equation}
 \frac{V_0\otimes V_1\otimes V_2}
 {W_0\otimes V_2+V_0\otimes W_1}
 \label{eq:intrinsic-cubic-envelope}
\end{equation}
である．この空間の次元は，対象間の同型やHom空間の基底の選択によらず，
アーベル圏同値で保たれる．
$F$の分類と補題\ref{lem:six-quadrics}によれば，その次元は十九である．
式\eqref{eq:orbit-product}の積の規約では，時間順に並べた合成は
反対環の積に対応するが，二次包絡の次元は変わらない．

他方，$\Phi^{-1}T^iF\simeq\mathcal B(i)$であり，
飽和性によって二次合成核は$G$の六本の二次関係と同一視される．
従って上の空間は$(G^{\mathrm{op}})_3$と同型で，
次元は二十でなければならない．
これは矛盾である．対象としてのSerre公式から合成図式だけを比較したため，
Nakayama捩れの選択やorbit環の積の同一視はこの結論に不要である．
\end{proof}

\begin{remark}
ここでHochschild homologyの非零な次数$\pm1$成分を，一般のAS三次超曲面に対する障害として使っていない．
障害は，任意の候補構造対象について二次関係の三次帰結がすでに四十五次元であり，
標準AS四次元ambientに必要な三次成分二十次元を確保できない，という有限次数の合成不変量である．
\end{remark}

\section{有限計算の証明書}
\label{sec:reproducibility}

判別曲線の滑らかさと二次包絡の次元は，次の有理数体上の計算で検証できる．
検証コードと入力は論文に付随する計算資料に含まれる．
\begin{enumerate}
\item \path{scripts/verify_degree3_smooth.py}は$q$の行列式を計算し，
四つのアフィンチャートで特異点イデアルのGr\"obner基底を求める．
四つの基底はいずれも$[1]$である．
\item \path{scripts/verify_degree3_relations.py}は
式\eqref{eq:six-quadrics}から$64\times48$整数行列を構成する．
\path{certificates/degree3_quadratic.json}に記載した小行列の行列式は$512$であり，
同ファイルの三本の整数右核ベクトルは独立である．
従って行列のrankは正確に四十五である．
このスクリプトは補題\ref{lem:numerical-rank-four}のEuler行列と，
命題\ref{prop:exceptional-rank-one}で用いる整数性の式も計算する．
\end{enumerate}
コードの生成元$(A,B,C,D)$は本文の$(x,y,z,w)$に対応する．
単項式の順序はこの順の辞書式順序とし，証明書の行・列の添字は零から数える．


\section{次数一・二のClifford反例と重み付き超曲面}
可換 del Pezzo 曲面の反標準環は，次数二では重み $(1,1,1,2)$ の四次式，
次数一では重み $(1,1,2,3)$ の六次式によって記述される．
非可換の場合には，反標準 Hilbert 級数と ambient の最小生成次数は
独立に区別すべきデータである．本稿はこの違いを，
構造対象を忘れたアーベル圏そのものの性質として示す明示例を与える．

線束値二次形式，偶 Clifford algebra，quaternion order と conic bundle の
対応は Chan--Ingalls~\cite{ChanIngalls2012} による既知の枠組みである．
特に本稿の次数一の例は，同論文の $F_{5-}^{2}$ 型の対称分解に属し，
滑らかな cubic threefold を直線に沿って blow up して得る conic bundle の
Clifford 圏でもある．この圏の数値的 Grothendieck 群については
Bernardara--Macr\`i--Mehrotra--Stellari~\cite{BMMS2012} を用いる．
従って，これらの order の型や Clifford 圏の存在自体を新しい発見として扱わない．

本稿で証明する中心的な事実は，具体的な二つの反標準環の表示と，
指定した最小生成次数の ambient に対する非表示である．
後者の核心は，任意のアーベル圏同値による自由対象の像を
局所射影な division rank 一の exceptional 対象に限定し，その全てを分類した上で，
隣接する三対象の合成写像の全射性を比較する点にある．
したがって，単に与えた構造対象に対する環の生成元を数えるだけではない．
存在する別の ambient は，graded Clifford algebra の正則性
\cite{CV2010,CKS2019} と Ore 拡大の正則性 \cite{CKZ2017,LWW2014} を用いて構成する．
この明示例と非表示定理の組合せについて，先行研究における初出は未確定である．

ここで「次数」$e$ は反標準 Hilbert 関数で正規化した次数である．
本稿の order は中心上で階数四，分岐指数二であり，
通常の order の標準因子は
\[
 K_{\mathcal A_e}=K_{Z_e}+\tfrac12D_e=-\tfrac12L_e,
 \qquad K_{\mathcal A_e}^{,2}=\tfrac e4.
\]
従って $e=\rk_{\mathcal O_{Z_e}}(\mathcal A_e)K_{\mathcal A_e}^{,2}$ である．
order の分類で用いる自己交点 \cite{Nasr2020} とこの正規化を混同しない．
Hochschild homology の計算には，quadric fibration の半直交分解
\cite{KuznetsovQuadric} と判別二重被覆による公式 \cite{KuznetsovHH} を用いる．


% !TEX root = ../nc-cubic-non-as.tex
\section{設定と主結果}
\label{sec:main}

基礎体を $k=\mathbb C$ とし，order の加群は右加群とする．
$B$ が連結局所有限・左右 noetherian 次数付き algebra のとき，
$\operatorname{qgr}B$ は有限生成次数付き右 $B$-module の圏を
有限次元加群で割った Serre 商圏であり，$\pi$ はその商関手を表す．
Hom と Ext が有限次元で高次 Ext が消滅する圏では
$\chi(E,F)=\sum_i(-1)^i\dim_k\operatorname{Ext}^i(E,F)$ と置き，
$K_{\mathrm{num}}$ は $K_0$ をこの Euler pairing の radical で割った群とする．
ここで使う圏は Serre functor を持つので，左右の radical は一致する．
導来圏の対象 $E$ が exceptional であるとは，
$\operatorname{Hom}(E,E)\simeq k$ かつ
$\operatorname{Hom}(E,E[i])=0$ for $i\ne0$ を満たすことをいう．

\begin{definition}
本稿で次数 $e$ の非可換 del Pezzo 曲面というとき，
noetherian な $k$-線形アーベル圏 $\mathcal C_e$ と対象 $\mathcal O_e$ に対し，
次の内在的条件を課す．圏の大域次元は二で，導来圏は smooth proper な
dg enhancement と Serre functor $S_e$ を持つ．
$T_e=S_e^{-1}[2]$ はアーベル圏の完全自己同値であり，
$\mathcal O_e$ は Gabriel--Krull rank 一の torsion-free exceptional 対象，
$(\mathcal O_e,T_e)$ は Artin--Zhang ample である．さらに
\[
 \dim\operatorname{Hom}(\mathcal O_e,T_e^n\mathcal O_e)
 =\begin{cases}1+\frac e2 n(n+1)&n\ge0,\\0&n<0,\end{cases}
 \qquad
 \operatorname{Ext}^1(\mathcal O_e,T_e^n\mathcal O_e)=0
 \quad(n\in\mathbb Z)
\]
を満たす．ここで torsion-free は低い Gabriel--Krull 次元の非零部分対象を
持たないという意味であり，order の例では低次元の台を持つ部分対象の不存在を表す．
\end{definition}

その反標準環を
\[
 R^{(e)}=\bigoplus_{n\ge0}
 \operatorname{Hom}_{\mathcal C_e}(\mathcal O_e,T_e^n\mathcal O_e)
\]
と定める．$f$ の次数が $m$，$g$ の次数が $n$ のとき，積の規約は
$fg=T_e^n(f)\circ g$ とする．

アーベル圏が曲面上の order の coherent module の圏である場合，
台の次元が $1$ 以下の対象で割った商圏における長さを
Gabriel--Krull rank と呼び，$\operatorname{GKrk}$ と書く．
これは台の次元そのものとは異なる．
以下では一般点の algebra が division algebra であることを証明するので，
構造対象の Gabriel--Krull rank は $1$ となる．

\begin{theorem}
\label{thm:weighted-main}
本稿で構成する $e=2,1$ の2つの組
$(\mathcal C_e,\mathcal O_e)$ は次の性質を持つ．
\begin{enumerate}
\item $\mathcal C_e$ の大域次元は二で，導来圏は smooth proper な
dg enhancement を持つ．$\mathcal O_e$ は torsion-free exceptional で，
$\operatorname{GKrk}(\mathcal O_e)=1$ である．
$(\mathcal O_e,T_e)$ は Artin--Zhang ample であり，
$\mathcal C_e\simeq\operatorname{qgr}R^{(e)}$ である．
\item すべての $n\ge0$ について
\[
 \dim_k R^{(e)}_n=1+\frac e2 n(n+1).
\]
従って
\[
 H_{R^{(2)}}(t)=\frac{1-t^4}{(1-t)^3(1-t^2)},\qquad
 H_{R^{(1)}}(t)=\frac{1-t^6}{(1-t)^2(1-t^2)(1-t^3)}.
\]
\item $e=2$ の場合，最小生成次数 $(1,1,1,2)$，
Hilbert 級数 $((1-t)^3(1-t^2))^{-1}$ の左右 noetherian
AS-regular algebra $G$ の次数4の正則正規元 $f$ に対して，
$\mathcal C_2\simeq\operatorname{qgr}(G/(f))$ となる
$k$-線形アーベル圏同値は存在しない．
$e=1$ の場合も，最小生成次数 $(1,1,2,3)$，
Hilbert 級数 $((1-t)^2(1-t^2)(1-t^3))^{-1}$，
次数6の正則正規元 $f$ に対して同じ非存在が成り立つ．
ここで $G$ の大域次元は $4$ とする．
\item 最小生成次数の条件を外すと，両例の反標準環は，
第3項でそれぞれの $G$ に指定した Hilbert 級数を持つ
大域次元 $4$ の左右 noetherian AS-regular algebra の中心正則超曲面になる．
実際の ambient の最小生成次数は，それぞれ $(1,1,1)$ と $(1,1,3)$ である．
\item 標準的な dg enhancement の Hochschild homology は
\[
\begin{array}{c|ccc}
e&\operatorname{HH}_{-1}(\mathcal C_e)&
\operatorname{HH}_0(\mathcal C_e)&\operatorname{HH}_1(\mathcal C_e)\\ \hline
2&k^3&k^4&k^3\\
1&k^5&k^3&k^5
\end{array}
\]
であり，$|i|\ge2$ では $\operatorname{HH}_i(\mathcal C_e)=0$ である．
\end{enumerate}
\end{theorem}

AS-regular は，有限大域次元，左右の AS-Gorenstein 条件，有限 GK 次元を
満たす連結局所有限 $\mathbb N$-次数付き algebra の意味で用いる．
大域次元を $d$ とすると，AS-Gorenstein parameter $\lambda$ は
$\underline{\operatorname{Ext}}^i_G(k,G)=0$ for $i\ne d$，
$\underline{\operatorname{Ext}}^d_G(k,G)\simeq k(\lambda)$
（および左加群での同じ条件）で定まる整数である．
最小生成次数は $G_+/(G_+)^2$ の次数付きベクトル空間としての次数を指す．
主定理の非表示は，構造対象を保つ同値だけでなく，
構造対象を指定しないアーベル圏同値も排除する．
一方，Hilbert 級数だけで ambient を指定する意味の非表示は主張しない．
第1，2項は本節と第\ref{sec:serre-rank}節，
第4項は第\ref{sec:explicit-ambient}節，
第3項は第\ref{sec:unpointed}節，
第5項は第\ref{sec:hochschild}節で証明する．

\subsection{Clifford algebra の規約}

$Z$ を滑らかな射影曲面，$L$ を線束，
$\mathcal V$ を階数 $3$ のベクトル束とし，
$q:\operatorname{Sym}^2\mathcal V\to L$ を対称形式とする．
局所的に $L$ と $\mathcal V$ を自明化したとき，Clifford relations は
$xy+yx=2q(x,y)$ とする．偶部分と奇部分は
\begin{equation}
 \mathcal A=\operatorname{Cl}_0(\mathcal V,q,L)
 =\mathcal O_Z\oplus(\wedge^2\mathcal V)\otimes L^{-1},
 \qquad
 \mathcal P=\operatorname{Cl}_1(\mathcal V,q,L)
 =\mathcal V\oplus(\det\mathcal V)\otimes L^{-1}
 \label{eq:clifford-bundles}
\end{equation}
である．これらの外積分解は，標数 $0$ における反対称化写像による
$\mathcal O_Z$-加群の同型であり，積は Clifford multiplication で定める．
$\mathcal P$ は $\mathcal A$-bimodule である．
$\mathbb P_Z(\mathcal V)$ は $\mathcal V$ の直線をパラメータ化する規約とし，
\[
 X=\{q=0\}\subset\mathbb P_Z(\mathcal V),\qquad p:X\to Z
\]
を付随する conic bundle とする．

\subsection{次数2の例}
\label{subsec:example2}

\[
 Z_2=\mathbb P^1_s\times\mathbb P^1_t,\qquad
 L_2=\mathcal O(1,1),\qquad \mathcal V_2=\mathcal O^{\oplus3}
\]
とする．Segre 座標
$a=s_0t_0,\ b=s_0t_1,\ c=s_1t_0,\ d=s_1t_1$ は $ad=bc$ を満たす．
二次形式を
\begin{equation}
 q_2=
 \begin{pmatrix}
 b&a+c&a+d\\
 a+c&b&b\\
 a+d&b&c
 \end{pmatrix}
 \label{eq:q2}
\end{equation}
で定める．その行列式は
\begin{align}
 \Delta_2={}&a^2b-a^2c+2abc-2ac^2-b^3+b^2c
             +2bcd-bd^2-c^3 .
 \label{eq:discriminant2}
\end{align}
$D_2=(\Delta_2=0)\subset Z_2$ は $(3,3)$ 曲線である．
式\eqref{eq:clifford-bundles}から
\begin{equation}
 \mathcal A_2=\mathcal O\oplus\mathcal O(-1,-1)^{\oplus3},
 \qquad
 \mathcal P_2=\mathcal O^{\oplus3}\oplus\mathcal O(-1,-1).
 \label{eq:bundles2}
\end{equation}
以後 $\mathcal C_2=\operatorname{coh}\mathcal A_2$，
$\mathcal O_2=\mathcal A_2$ とする．

\subsection{次数1の例}
\label{subsec:example1}

\[
 Z_1=\mathbb P^2_{[u:v:w]},\qquad
 L_1=\mathcal O(1),\qquad
 \mathcal V_1=\mathcal O^{\oplus2}\oplus\mathcal O(-1)
\]
とし，
\begin{equation}
 q_1=
 \begin{pmatrix}
 u&v&w^2\\
 v&w&u^2\\
 w^2&u^2&2v^3
 \end{pmatrix}
 \label{eq:q1}
\end{equation}
で定める．その行列式は
\begin{equation}
 \Delta_1=-u^5+2u^2vw^2+2uv^3w-2v^5-w^5.
 \label{eq:discriminant1}
\end{equation}
$D_1=(\Delta_1=0)\subset\mathbb P^2$ は5次曲線である．このとき
\begin{equation}
 \begin{split}
 \mathcal A_1&=\mathcal O\oplus\mathcal O(-1)\oplus\mathcal O(-2)^{\oplus2},\\
 \mathcal P_1&=\mathcal O^{\oplus2}\oplus\mathcal O(-1)\oplus\mathcal O(-2).
 \end{split}
 \label{eq:bundles1}
\end{equation}
$\mathcal C_1=\operatorname{coh}\mathcal A_1$，
$\mathcal O_1=\mathcal A_1$ とする．

\begin{lemma}
\label{lem:smooth}
$D_2,D_1$ は滑らかで連結であり，それぞれ種数 $4,6$ を持つ．
また $q_e$ の rank はすべての点で $2$ 以上で，
$p_e:X_e\to Z_e$ は滑らかな全空間を持つ flat conic bundle である．
\end{lemma}
\begin{proof}
滑らかさは，付録\ref{sec:computation}で指定する厳密計算による．
各標準アフィンチャートで，判別式 $F$ と2個の偏微分が生成する
$\mathbb Q$ 上のイデアルの Gröbner 基底は $\{1\}$ である．
$Z_2$ の4チャートと $Z_1$ の3チャートが全体を被覆するので，
$D_e$ は $\mathbb C$ 上で滑らかである．
完全列 $0\to\mathcal O_{Z_e}(-D_e)\to\mathcal O_{Z_e}\to
\mathcal O_{D_e}\to0$ について，$H^1(\mathcal O_{Z_e}(-D_e))=0$ なので
$H^0(\mathcal O_{D_e})=k$ であり，連結である．
随伴公式は $g(D_2)=(3-1)(3-1)=4$，
$g(D_1)=(5-1)(5-2)/2=6$ を与える．

対称行列の rank が $1$ 以下なら全ての $2\times2$ minor が消え，
行列式の微分も消えるため，$D_e$ の滑らかさに反する．
従って各ファイバーは非零な2次式で定義され，
一定の Hilbert 多項式を持つので $p_e$ は flat である．
$D_e$ 外ではファイバーが滑らかである．$D_e$ 上では rank が $2$ で，
退化 conic の特異点における基底方向の微分は
判別式の非零微分に対応する．従って全空間 $X_e$ は滑らかである．
\end{proof}

\section{一般点の rank と反標準自己同値}
\label{sec:serre-rank}

この節では，添字 $e$ を省略する．両例に共通する等式は
\begin{equation}
 \det\mathcal V=\omega_Z\otimes L^2,\qquad L\ \text{は ample}.
 \label{eq:detV}
\end{equation}

\subsection{判別二重被覆と division 性}
\label{subsec:residue}

$i:D\hookrightarrow Z$ とし，$q:\mathcal V\to\mathcal V^\vee\otimes L$ の
余核を $i_*\mathcal M$ と書く．$\det q$ が非零で，$D$ 上の corank が
ちょうど $1$ なので $\mathcal M$ は $D$ 上の線束である．
$q|_D$ の radical を $\mathcal K$，
$\mathcal W=\mathcal V|_D/\mathcal K$ とすると，$\mathcal W$ は
$L|_D$ に値を持つ非退化な階数2の二次束であり，
\[
 \mathcal M=\mathcal K^\vee\otimes L|_D,\qquad
 \eta:=\det\mathcal W\otimes L^{-1}|_D
 =\mathcal M\otimes(\det\mathcal V)|_D\otimes L^{-2}|_D.
\]
誘導された非退化形式を $q_{\mathcal W}$ とし，その determinant の負号を取った同型
\[
 \vartheta=-\det(q_{\mathcal W}):\eta^{\otimes2}\xrightarrow{\sim}\mathcal O_D
\]
を用いて，$\mathcal O_D\oplus\eta$ に
$(f,x)(g,y)=(fg+\vartheta(x\otimes y),fy+gx)$ という積を入れる．
これにより有限非分岐二重被覆
\[
 \pi:\widetilde D=\operatorname{Spec}_D(\mathcal O_D\oplus\eta)\longrightarrow D,
 \qquad
 \pi_*\mathcal O_{\widetilde D}=\mathcal O_D\oplus\eta
\]
を得る．局所的に $L|_D$ を自明化し，$q_{\mathcal W}=\langle a,b\rangle$ と書くと，
$\eta$ の外積基底に対応する元 $\xi$ は $\xi^2=-ab$ を満たす．
等方直線 $[x:y]$ の式は $ax^2+by^2=0$ であり，$\xi=by/x$ によって
この二重被覆の2点と対応する．等方直線の族
$\rho:F_D=\{q_{\mathcal W}=0\}\subset\mathbb P_D(\mathcal W)\to D$ と，
その tautological line $\mathcal U\subset\rho^*\mathcal W$ を取る．この対応は
\[
 \rho^*\det\mathcal W=\mathcal U\otimes(\rho^*\mathcal W/\mathcal U)
 \xrightarrow{q_{\mathcal W}}\rho^*(L|_D)
\]
から定まるので局所基底に依存せず，$F_D\simeq\widetilde D$ を与える．
各等方直線の $\mathcal V|_D$ 内の逆像を射影化すると退化 conic の1成分になるので，
$\widetilde D$ はその2つの成分を区別する二重被覆である．
式\eqref{eq:detV}より $\eta=\mathcal M\otimes\omega_Z|_D$ なので
\begin{equation}
 0\longrightarrow \mathcal V\otimes\omega_Z
 \xrightarrow{q} \mathcal V^\vee\otimes L\otimes\omega_Z
 \longrightarrow i_*\eta\longrightarrow0 .
 \label{eq:eta-resolution}
\end{equation}
次数2では左の2項は
$\mathcal O(-2,-2)^3,\mathcal O(-1,-1)^3$，
次数1では
$\mathcal O(-3)^2\oplus\mathcal O(-4),
\mathcal O(-2)^2\oplus\mathcal O(-1)$ である．
前者の $H^1$ と後者の $H^0$ はともに消滅する．従って
\begin{equation}
 H^0(D,\eta)=0,\qquad \eta\not\simeq\mathcal O_D.
 \label{eq:eta-nontrivial}
\end{equation}
$D$ が連結なので，$\widetilde D$ も連結である．

\begin{proposition}
\label{prop:division}
$\mathscr D=\mathcal A\otimes_{\mathcal O_Z}k(Z)$ は quaternion division algebra
であり，$\operatorname{GKrk}(\mathcal A)=1$ である．
\end{proposition}
\begin{proof}
非退化な3変数二次形式の偶 Clifford algebra は quaternion algebra である．
$D$ の一般点の離散付値環上で形式を
$\langle a,b,sc\rangle$ と対角化する．ここで $a,b,c$ は単元，
$s$ は一様化元である．偶 Clifford algebra の生成元
$I=e_1e_2,\ J=e_1e_3$ は
$I^2=-ab,\ J^2=-asc,\ IJ=-JI$ を満たす．
その residue は $\sqrt{-ab}$ を加える二次拡大であり，
これは $\widetilde D\to D$ の関数体拡大に等しい．
式\eqref{eq:eta-nontrivial}よりこの residue は非自明なので，
$\mathscr D$ の Brauer 類は非零である．次数2の central simple algebra
は split でなければ division だから，$\mathscr D$ は division である．

$\mathcal C_{\le1}\subset\operatorname{coh}\mathcal A$ を，台の次元が $1$ 以下の対象の
Serre 部分圏とする．一般点への制限は
$\operatorname{coh}\mathcal A/\mathcal C_{\le1}\simeq\operatorname{mod}\mathscr D$
を与える．有限次元 $\mathscr D$-module とその射は分母を払って
稠密開集合上に延長でき，2つの延長の差は一般点で消える対象で測られるためである．
正則右 $\mathscr D$-module の長さは $1$ なので結論を得る．
\end{proof}

以後，一般点の右 $\mathscr D$-module が $\mathscr D^{\oplus r}$ であるとき，
その整数 $r$ を division rank と呼ぶ．

\subsection{局所正則性と Serre functor}

閉点の完備局所環を $R=k[[s,t]]$ とする．$D$ 外では，非退化形式を対角化し，
単元の平方根を取ることにより $\widehat{\mathcal A}\simeq M_2(R)$ を得る．
$D=(s)$ 上では，$L$ と $\mathcal V$ を自明化して
$q=\langle a,b,sc\rangle$ と対角化できる．ここで $a,b,c\in R^\times$ であり，
最後の係数の位数が $1$ であることは $D$ の滑らかさによる．
$R$ の単元は平方根を持つので，対応する Clifford の偶元を正規化して
\[
 I=\frac{e_1e_2}{\sqrt{-ab}},\qquad
 J=\frac{e_1e_3}{\sqrt{-ac}},\qquad
 I^2=1,\quad J^2=s,\quad IJ=-JI
\]
とできる．$R$-基底 $1,I,J,IJ$ に対し
\[
 I\longmapsto\begin{pmatrix}1&0\\0&-1\end{pmatrix},\qquad
 J\longmapsto\begin{pmatrix}0&1\\s&0\end{pmatrix}
\]
は積を保ち，その像は次の order 全体を生成する．従って
\begin{equation}
 \widehat{\mathcal A}\simeq
 \begin{pmatrix}R&R\\sR&R\end{pmatrix}.
 \label{eq:local-order}
\end{equation}
この order の大域次元を確認するため，$U=k[[r,t]]$，$r^2=s$ とし，
$g(r)=-r,\ g(t)=t$ で作用する位数 $2$ の群を用いて
$\Lambda=U\rtimes\langle g\rangle$ と置く．その積は
$g^2=1,\ gu=g(u)g$（$u\in U$）で定める．$U=Rr\oplus R$ 上の作用では，
$r$ と $g$ の行列はそれぞれ上の $J$ と $-I$ である．
従って $R$-基底 $1,g,r,rg$ の像を比較すると
$\Lambda\simeq\widehat{\mathcal A}$ となる．
任意の右 $\Lambda$-module $M$ に対し，積による全射
$M\otimes_U\Lambda\to M$ は
\[
 m\longmapsto\frac12\bigl(m\otimes1+mg\otimes g\bigr)
\]
で右 $\Lambda$-線形に分裂する．$U$ は大域次元 $2$ の正則局所環で，
$\Lambda$ は左右の $U$-module として自由なので，$U$ 上の長さ $2$ 以下の射影分解を
$-\otimes_U\Lambda$ で誘導すれば $\operatorname{pd}_{\Lambda}M\le2$ となる．
また $r=t=0,\ g=1$ で作用する加群 $k$ は，$U$ 上の射影次元が $2$ であり，
$\Lambda$ 上の射影分解を $U$ に制限できることから
$\operatorname{pd}_{\Lambda}k\ge2$ である．従って $\operatorname{gl.dim}\Lambda=2$ である．
分岐外でも $M_2(R)$ の大域次元は $2$ である．
有限表示加群の射影性は中心の忠実平坦な基底変換で降下するので，
第2 syzygy にこれを適用すると，元の局所 order 上でも coherent module は
長さ $2$ 以下の射影分解を持つ．

各点の近傍で $q$ は単元を表すので，対応する Clifford の奇元は可逆である．
その奇元による局所自明化を使うと，Clifford multiplication は
\begin{equation}
 \mathcal P\otimes_{\mathcal A}\mathcal P\simeq\mathcal A\otimes L
 \label{eq:P-square}
\end{equation}
となる．射は元の multiplication から定まり，局所的に同型であるため
選択に依存せず大域的な同型である．
同じく最高外冪成分への pairing
$\mathcal A\otimes\mathcal P\to\det\mathcal V\otimes L^{-1}$
は完全であり，偶元と奇元について cyclic である．従って bimodule として
\[
 \mathcal A^\vee\simeq\mathcal P\otimes(\det\mathcal V)^{-1}\otimes L.
\]
$\omega_{\mathcal A}:=\mathcal Hom_{\mathcal O_Z}(\mathcal A,\omega_Z)$ と置くと，
式\eqref{eq:detV}，\eqref{eq:P-square}は
\begin{equation}
 \omega_{\mathcal A}\simeq\mathcal P\otimes L^{-1}\simeq\mathcal P^{-1}
 \label{eq:dualizing-order}
\end{equation}
を与える．$Z$ 上の Grothendieck--Serre duality を
局所射影分解に適用すると
\begin{equation}
 S=(-\otimes_{\mathcal A}\mathcal P^{-1})[2],\qquad
 T=S^{-1}[2]=-\otimes_{\mathcal A}\mathcal P,\qquad
 T^2=-\otimes L.
 \label{eq:Serre}
\end{equation}
特に $S[-2]$ と $T$ はアーベル圏の完全自己同値である．

\subsection{構造対象についての AZ ampleness}

任意の coherent $\mathcal A$-module $M$ は coherent $\mathcal O_Z$-module
でもある．$m\gg0$ に対して $M\otimes L^m$ が大域生成されるので，
随伴により $\mathcal A^{\oplus N}\twoheadrightarrow M\otimes L^m$ が得られる．
従って $(T^{-2m}\mathcal A)^{\oplus N}\twoheadrightarrow M$ となり，
ample pair の生成条件を満たす．

全射 $M\twoheadrightarrow N$ の核を $K$ とし，$n=2m+r$，
$r\in\{0,1\}$ と書く．$\mathcal A$ は局所射影なので
\[
 \operatorname{Ext}^1_{\mathcal A}(\mathcal A,T^nK)
 =H^1\bigl(Z,(K\otimes_{\mathcal A}\mathcal P^r)\otimes L^m\bigr)=0
 \qquad(m\gg0).
\]
これは ample pair の射の全射性の条件である．従って
$(\mathcal A,T)$ は Artin--Zhang ample である．
構造対象が noetherian，その自己準同型環が有限次元，
すべての Hom が有限次元という条件も成り立つので，
Artin--Zhang の再構成定理 \cite[Theorem~4.5(1)]{AZ1994} により
\begin{equation}
 \operatorname{coh}\mathcal A\simeq\operatorname{qgr}R,\qquad
 R_n=H^0(Z,\mathcal P^{\otimes_{\mathcal A}n})
 \label{eq:AZ-reconstruction}
\end{equation}
を得る．

\subsection{Hilbert 級数と exceptionality}

式\eqref{eq:Serre}から
$R_{2m}=H^0(Z,\mathcal A\otimes L^m)$，
$R_{2m+1}=H^0(Z,\mathcal P\otimes L^m)$ である．
式\eqref{eq:bundles2}，\eqref{eq:bundles1}に代入すると
\[
\begin{array}{c|cc}
 e&\dim R^{(e)}_{2m}&\dim R^{(e)}_{2m+1}\\ \hline
2&(m+1)^2+3m^2=4m^2+2m+1&
3(m+1)^2+m^2=4m^2+6m+3\\
1&2m^2+m+1&2m^2+3m+2
\end{array}
\]
となる．次数1の行では
$h^0(\mathbb P^2,\mathcal O(r))=\binom{r+2}{2}$ を用いるが，
現れる $r\ge-2$ について $r=-1,-2$ の値は $0$ である．
従ってこの表は $m=0$ を含むすべての $m\ge0$ で成り立つ．
これより主定理の Hilbert 関数と Hilbert 級数が従う．
また
\[
 \operatorname{Ext}^i_{\mathcal A}(\mathcal A,\mathcal A)=H^i(Z,\mathcal A)
 =\begin{cases}k&i=0,\\0&i\ne0,\end{cases}
\]
なので $\mathcal A$ は exceptional である．

\begin{proposition}
\label{prop:intrinsic-completion}
上の二つの組は，定義に挙げた残りの内在的条件も満たす．
\end{proposition}
\begin{proof}
$\mathcal A$ は中心上のベクトル束なので，低次元の台を持つ非零部分加群を持たず，
torsion-free である．$n<0$ では
$\mathcal A\otimes L^m$ ($n=2m$) と
$\mathcal P\otimes L^m$ ($n=2m+1$) の各直和成分は負の線束であり，
大域切断を持たない．次数二では全て $\mathcal O(r,r)$ の形で，
K\"unneth 公式より $H^1(\mathcal O(r,r))=0$ が全整数 $r$ で成り立つ．
次数一では $H^1(\mathbb P^2,\mathcal O(r))=0$ が全整数 $r$ で成り立つ．
従ってすべての $n$ で所定の $\operatorname{Ext}^1$ は消滅する．

式\eqref{eq:Serre}とアーベル圏の負の Ext の消滅から，
任意の coherent $M,N$ と $i>2$ について
\[
 \operatorname{Ext}^i(M,N)
 \simeq\operatorname{Hom}\bigl(N,M\otimes_{\mathcal A}\mathcal P^{-1}[2-i]\bigr)^*=0.
\]
一方 $\operatorname{Ext}^2(\mathcal A,\omega_{\mathcal A})
\simeq\operatorname{End}(\mathcal A)^*\ne0$ なので大域次元は二である．
局所射影分解により $D^b(\operatorname{coh}\mathcal A)$ は
$\operatorname{Perf}(Z,\mathcal A)$ と一致する．
滑らかな射影三次元多様体 $X$ に対する Kuznetsov の半直交分解
\cite[Theorem~4.2]{KuznetsovQuadric} はこの圏を admissible 成分として実現する．
対応する dg 圏は smooth proper な $\operatorname{Perf}_{\mathrm{dg}}(X)$ の
admissible 成分なので，smooth proper である．
\end{proof}


% !TEX root = ../nc-cubic-non-as.tex
\section{反標準環の明示表示と AS-regular ambient}
\label{sec:explicit-ambient}

Clifford 積を用いて反標準環の次数付き algebra としての表示を与え，
その ambient が AS-regular であることを示す．

以下で用いる Clifford の PBW 性を確認する．
$S$ を可換 $\mathbb C$-algebra，$Q=(q_{ij})$ を $S$ 上の対称行列とし，
\[
 C_S(Q)=S\langle x_1,\ldots,x_r\rangle/
 (x_ix_j+x_jx_i-2q_{ij})
\]
と置く．$S$ はこの表示で中心に作用するものとする．
順序付き積 $x_{i_1}\cdots x_{i_s}$，$i_1<\cdots<i_s$ は $S$-基底である．
実際，関係式によりこれらは生成系となる．
独立性には $E=S^r$ 上の外積代数を用い，$x_i$ を
$c_i=e_i\wedge(-)+\iota_{q(e_i,-)}$ として作用させる．
$c_ic_j+c_jc_i=2q_{ij}$ であり，$c_{i_1}\cdots c_{i_s}(1)$ の最高外積次数の成分は
$e_{i_1}\wedge\cdots\wedge e_{i_s}$ である．
従ってこの作用を $1$ に適用し，外積次数について降下すれば独立性が従う．
この議論は $Q$ の非退化性を仮定しない．

\subsection{次数2の反標準環}

$S=\mathbb C[a,b,c,d]$，$\deg a=\deg b=\deg c=\deg d=2$ とし，
\[
 Q_2(a,b,c,d)=
 \begin{pmatrix}
 b&a+c&a+d\\
 a+c&b&b\\
 a+d&b&c
 \end{pmatrix}.
\]
$\deg X_i=1$ として
\begin{equation}
 G^{(2)}=C_S(Q_2)
 \label{eq:ambient-d2}
\end{equation}
を定める．PBW 性から，$G^{(2)}$ は $S$ 上，
\[
 1,\quad X_1,X_2,X_3,\quad X_1X_2,X_1X_3,X_2X_3,
 \quad X_1X_2X_3
\]
を基底とする自由加群である．従って
\begin{equation}
 H_{G^{(2)}}(t)
 =\frac{(1+t)^3}{(1-t^2)^4}
 =\frac{1}{(1-t)^3(1-t^2)}.
 \label{eq:ambient-d2-hilbert}
\end{equation}

中心元 $h=ad-bc$ は $S$ 上で非零因子なので，$G^{(2)}$ 上でも正則である．
$\mathcal V_2=\mathcal O^{\oplus3}$ の標準基底の切断を
$\varepsilon_1,\varepsilon_2,\varepsilon_3$ とし，これらを
\[
 H^0(Z_2,\mathcal V_2)\subset H^0(Z_2,\mathcal P_2)=R^{(2)}_1
\]
の元と見る．また，Segre 切断 $a,b,c,d$ をスカラー成分
\[
 H^0(Z_2,L_2)\subset H^0(Z_2,\mathcal A_2\otimes L_2)=R^{(2)}_2
\]
の元と見る．これらの次数2の切断は Clifford 積に関して中心にあり，
$\varepsilon_i\varepsilon_j+\varepsilon_j\varepsilon_i=2(Q_2)_{ij}$ と
$ad-bc=0$ を満たす．従って $X_i\mapsto\varepsilon_i$ と，
$a,b,c,d$ をそれぞれ同名の Segre 切断へ送る対応により，次数付き algebra の射
\begin{equation}
 G^{(2)}/(ad-bc)\longrightarrow R^{(2)}
 \label{eq:d2-hypersurface-map}
\end{equation}
を定める．これが全次数で同型であることを確認する．
Segre 埋め込みの切断環を，次数を $2$ 倍して
\[
 S'=\bigoplus_{m\geq0}H^0(\mathbb P^1\times\mathbb P^1,
 \mathcal O(m,m)),\qquad \deg H^0(\mathcal O(m,m))=2m
\]
と書けば，$S'=S/(ad-bc)$ である．
Clifford の外積分解から $R^{(2)}$ は $S'$ 上自由で，その基底の次数は
\[
 0,\quad 1,1,1,\quad 2,2,2,\quad 3
\]
である．上記の順序付き Clifford 積は，その外積分解に対して対角成分が $1$ の
三角形の基底変換を与える．従って \eqref{eq:d2-hypersurface-map} は $S'$-加群として
同型であり，algebra としても同型である．

\begin{proposition}
$G^{(2)}$ は大域次元 $4$ の左右 noetherian AS-regular algebra であり，
最小生成元の次数は $(1,1,1)$ である．
従って $R^{(2)}$ は Hilbert 級数
$((1-t)^3(1-t^2))^{-1}$ を持つ AS-regular algebra の中心正則 $4$ 次超曲面である．
\end{proposition}

\begin{proof}
行列成分から
\[
 b=X_1^2,\qquad c=X_3^2,\qquad
 a=\frac{X_1X_2+X_2X_1}{2}-c,\qquad
 d=\frac{X_1X_3+X_3X_1}{2}-a
\]
を得る．従って $G^{(2)}$ は $X_1,X_2,X_3$ で生成され，これらは次数 $1$ の
独立な元なので最小生成系である．

PBW 性から $d$ は中心正則元である．対称係数行列を
\[
 M_a=\begin{pmatrix}0&1&1\\1&0&0\\1&0&0\end{pmatrix},\qquad
 M_b=\begin{pmatrix}1&0&0\\0&1&1\\0&1&0\end{pmatrix},\qquad
 M_c=\begin{pmatrix}0&1&0\\1&0&0\\0&0&1\end{pmatrix}
\]
と置けば，$Q_2|_{d=0}=aM_a+bM_b+cM_c$ である．
中心生成元を $2a,2b,2c$ と取ると，$G^{(2)}/(d)$ は係数行列
$M_a,M_b,M_c$ の graded Clifford algebra となる．
可換変数の列ベクトル $\boldsymbol{\xi}=(x,y,z)^{\mathsf t}$ に対し，
対応する二次形式は $\boldsymbol{\xi}^{\mathsf t}M_a\boldsymbol{\xi}$，
$\boldsymbol{\xi}^{\mathsf t}M_b\boldsymbol{\xi}$，
$\boldsymbol{\xi}^{\mathsf t}M_c\boldsymbol{\xi}$，すなわち
\[
 2x(y+z),\qquad x^2+y^2+2yz,\qquad 2xy+z^2
\]
である．Basepoint-free criterion \cite[Theorem~1.1]{CV2010} によれば，これらの二次形式が
$\mathbb P^2$ に共通零点を持たなければ，$G^{(2)}/(d)$ は大域次元 $3$ の
AS-regular algebra である．
実際，第 $1$ 式から $x=0$ または $y=-z$ である．
$x=0$ なら第 $3$ 式，第 $2$ 式の順に $z=y=0$ となる．
$x\neq0$ かつ $y=-z$ なら第 $3$ 式は $z(z-2x)=0$ となり，
$z=0$ と $z=2x$ の場合に第 $2$ 式はそれぞれ $x^2$ と $-3x^2$ になる．
従って共通の射影零点はなく，この criterion を適用できる．
正則正規元による拡大の定理 \cite[Proposition~2.1]{CKS2019} を $d$ に適用すると，
$G^{(2)}$ は大域次元 $4$ の AS-regular algebra となる．
Noetherian 性は，中心の多項式環 $S$ 上の有限加群であることから従い，
有限 GK 次元は \eqref{eq:ambient-d2-hilbert} から従う．
\end{proof}

\subsection{次数1の反標準環}

\begin{equation}
 H=\mathbb C\langle X,Y\rangle/
 (X^2Y-YX^2,\;Y^2X-XY^2),\qquad \deg X=\deg Y=1
 \label{eq:down-up-base}
\end{equation}
と置く．これは down-up algebra $D(0,1)$ である．
従って大域次元 $3$ の左右 noetherian AS-regular domain であり，
\[
 H_H(t)=\frac1{(1-t)^2(1-t^2)}
\]
となる \cite[Lemma~1.1]{CKZ2017}．

$u=X^2$，$w=Y^2$，$v=(XY+YX)/2$ は $H$ の中心元である．
特に $v$ の中心性は
\[
 [X,XY+YX]=X^2Y-YX^2=0,\qquad
 [Y,XY+YX]=Y^2X-XY^2=0
\]
から従う．
形式的な多項式環 $S_0=\mathbb C[u,v,w]$ 上の Clifford algebra
\[
 S_0\langle X,Y\rangle/(X^2-u,\;Y^2-w,\;XY+YX-2v)
\]
と $H$ の間には，これらの式による互いに逆な準同型がある．
従って PBW 性から $u,v,w$ は代数的に独立であり，
$H$ は $S_0$ 上，$1,X,Y,XY$ を基底とする自由加群である．

自己同型 $\sigma$ と $\sigma$-derivation $\delta$ を
\[
 \sigma(X)=-X,\quad\sigma(Y)=-Y,\qquad
 \delta(X)=2Y^4,\quad\delta(Y)=2X^4
\]
で定める．ここでは $\delta(ab)=\sigma(a)\delta(b)+\delta(a)b$ を規約とする．
$X^2,Y^2$ の中心性を使うと $\delta(X^2)=\delta(Y^2)=0$ であり，
\[
 \delta(XY+YX)
 =-2X^5+2Y^5-2Y^5+2X^5=0.
\]
従って $\delta(u)=\delta(v)=\delta(w)=0$ である．
また \eqref{eq:down-up-base} の関係に $\delta$ を適用した結果は
それぞれ $[X^2,2X^4]=0$，$[Y^2,2Y^4]=0$ なので，$\delta$ は $H$ 上に降りる．

\begin{equation}
 G^{(1)}=H[Z;\sigma,\delta],\qquad
 Z a=\sigma(a)Z+\delta(a),\qquad \deg Z=3
 \label{eq:ambient-d1}
\end{equation}
と置く．$\sigma$ は次数を保ち，$\delta$ は次数 $3$ を持つので，
これは連結次数付き Ore extension である．
$\sigma\delta+\delta\sigma=0$ は $X,Y$ 上で成立し，その積公式により $H$ 全体で成立する．
そのため $\delta^2$ は通常の derivation となる．
$\delta^2(X)=\delta^2(Y)=0$ なので $\delta^2=0$ である．
よって任意の $a\in H$ に対して
\[
 Z^2a=\sigma^2(a)Z^2+
 (\delta\sigma+\sigma\delta)(a)Z+\delta^2(a)=aZ^2.
\]
従って $Z^2$ は中心元であり，$u,v,w$ も $G^{(1)}$ の中心元である．

\begin{proposition}
$G^{(1)}$ は大域次元 $4$ の左右 noetherian AS-regular domain であり，
\[
 H_{G^{(1)}}(t)=\frac1{(1-t)^2(1-t^2)(1-t^3)}.
\]
最小生成元の次数は $(1,1,3)$ であり，次数 $6$ の中心正則元について
\begin{equation}
 R^{(1)}\simeq G^{(1)}/(Z^2-2v^3)
 \label{eq:d1-hypersurface}
\end{equation}
が成り立つ．
\end{proposition}

\begin{proof}
自己同型を伴う Ore extension は noetherian 性と domain 性を保つ．
また \cite[Theorem~3.3 and Remark~2.2]{LWW2014} を用いると，
正次数の次数付き Ore extension は twisted Calabi--Yau 次元を $1$ 増やす．
連結次数付きの場合の AS-regular 性との同値
\cite[Proposition~1.6]{CKS2019}により，$G^{(1)}$ は
大域次元 $4$ の AS-regular algebra である．
Ore の自由加群分解 $G^{(1)}=\bigoplus_{n\geq0}HZ^n$ は上記の Hilbert 級数を与える．
また，この分解から $Z$ は $X,Y$ で生成される部分 algebra に属さないので，
$X,Y,Z$ は次数 $(1,1,3)$ の最小生成系である．

$f=Z^2-2v^3$ は中心元である．Ore の自由加群分解における $Z^2$ の係数が $1$ なので
$f\neq0$ であり，domain 性から正則である．
$Z$ について monic な関係 $f=0$ で除算すれば，$G^{(1)}/(f)$ は
$H$ 上 $1,Z$ を基底とする自由加群となる．
従って $S_0$ 上の基底とその次数は
\[
 \begin{array}{c|cccccccc}
 \text{基底}&1&X&Y&XY&Z&XZ&YZ&XYZ\\\hline
 \text{次数}&0&1&1&2&3&4&4&5
 \end{array}
\]
である．定義関係は
\[
 X^2=u,\quad Y^2=w,\quad XY+YX=2v,\quad
 XZ+ZX=2w^2,\quad YZ+ZY=2u^2,\quad Z^2=2v^3.
\]
これらは次数 $1$ の例の二次形式
\[
 \begin{pmatrix}
 u&v&w^2\\ v&w&u^2\\ w^2&u^2&2v^3
 \end{pmatrix}
\]
の Clifford relations である．
ここで $X,Y$ は $H^0(\mathcal P_1)$，$Z$ は
$H^0(\mathcal P_1\otimes\mathcal O(1))$ の $\mathcal O$ 成分から取る．
Clifford multiplication により $G^{(1)}/(f)\to R^{(1)}$ が定まる．
一方，$R^{(1)}$ は $S_0=\bigoplus_{m\geq0}H^0(\mathbb P^2,\mathcal O(m))$
（切断次数を $2$ 倍する）上，同じ $8$ 個の次数を持つ自由加群であり，
順序付き Clifford 積の外積成分は上記の PBW 基底を与える．
従ってこの準同型は全次数で同型である．
\end{proof}

\subsection{次数2の積}

\begin{lemma}
\label{lem:mu2-surjective}
両方の反標準環で，積写像
\[
 \mu^{(e)}_2:R^{(e)}_1\otimes_{\mathbb C}R^{(e)}_1\longrightarrow R^{(e)}_2
 \qquad(e=1,2)
\]
は全射である．
\end{lemma}
\begin{proof}
次数2の例では，$R^{(2)}$ は次数1の $X_1,X_2,X_3$ の像で生成される．
次数1の例では，$\deg Z=3$ と $\deg(Z^2-2v^3)=6$ より，
式\eqref{eq:d1-hypersurface}を通じて $R^{(1)}_1=H_1$，$R^{(1)}_2=H_2$ と同定される．
$H$ は次数1の $X,Y$ で生成されるので，$H_2=H_1H_1$ である．
\end{proof}


% !TEX root = ../nc-cubic-non-as.tex
\section{構造対象を忘れた超曲面表示の排除}
\label{sec:unpointed}

この節では，最小生成次数を指定した主定理の非表示を証明する．
最初に局所射影な rank 1 の exceptional object を分類し，
その後に仮定したアーベル圏同値の下で構造対象の像が
この分類に入ることを示す．

\subsection{局所正則型と数値的類}

\begin{lemma}
\label{lem:regular-type}
$F$ が $\operatorname{coh}\mathcal A_e$ の局所射影対象で，
一般点での division rank が $1$ なら，
$F$ は局所的に正則右加群 $\mathcal A_e$ と同型である．
特にある $N\in\operatorname{Pic}(Z_e)$ に対し
\begin{equation}
 \det F=\det\mathcal A_e\otimes N^2
 \label{eq:norm-determinant}
\end{equation}
となる．
\end{lemma}
\begin{proof}
$D_e$ の一般点の離散付値環上では，
命題\ref{prop:division}の表示の $R[I]$ は非自明な剰余二次拡大を持つ．
$\mathcal A_e=R[I]\oplus R[I]J$，$J^2$ は一様化元の単元倍であり，
$J\mathcal A_e$ が Jacobson radical，その商は二次剰余体である．
従って $\mathcal A_e$ はこの点で局所環であり，
有限射影加群 $F$ は自由である．一般点での rank から自由階数は $1$ となる．

閉分岐点で完備化し，式\eqref{eq:local-order}の
2つの原始冪等元に対応する射影加群を $P_1,P_2$ とする．
完備環は semiperfect なので
$\widehat F=P_1^{\oplus r_1}\oplus P_2^{\oplus r_2}$ と書ける．
形式的な分岐因子 $(s)$ の一般点で局所化すると，直前の自由性が
基底変換される．そこでの2つの射影型は異なるため
$r_1=r_2=1$ であり，$\widehat F\simeq\widehat{\mathcal A_e}$ である．
分岐外では $\widehat{\mathcal A_e}=M_2(R)$ であり，
同じ $\mathcal O_Z$-rank $4$ の射影加群は正則加群と同型になる．
完備化上の生成元を中心の極大イデアル modulo で持ち上げ，
Nakayama の補題を適用すると，元の局所環でも
$\mathcal A_e\to F$ は全射になる．同じ階数のベクトル束の全射なので同型である．

従って $F$ の遷移写像は $\mathcal A_e^\times$ の元 $u$ による左乗法である．
quaternion algebra 上の左乗法の determinant は
$\det(L_u)=\operatorname{Nrd}(u)^2$ である．
reduced norm は乗法的で，単元を $\mathcal O_Z^\times$ に写すので，
遷移関数 $\operatorname{Nrd}(u)$ が線束 $N$ を定め，
式\eqref{eq:norm-determinant}を得る．
\end{proof}

\begin{lemma}
\label{lem:numerical-to-isomorphism}
$F$ を局所射影な division rank 1 の対象とし，
$[F]=[\mathcal A_e]$ in $K_{\mathrm{num}}(\mathcal A_e)$ とする．
このとき $F\simeq\mathcal A_e$ である．
\end{lemma}
\begin{proof}
忘却関手 $\operatorname{Forg}:\operatorname{coh}\mathcal A_e
\to\operatorname{coh}Z_e$ の左随伴は
$E\mapsto E\otimes_{\mathcal O_{Z_e}}\mathcal A_e$ なので，
\[
 \chi_{Z_e}(E,\operatorname{Forg}F)
 =\chi_{\mathcal A_e}(E\otimes\mathcal A_e,F).
\]
従って忘却関手は数値的 Grothendieck 群上に降り，
$c_1(F)=c_1(\mathcal A_e)$ となる．$Z_e$ の Picard 群は torsion-free で
$\operatorname{Pic}^0=0$ だから $\det F\simeq\det\mathcal A_e$ である．
$H=c_1(\mathcal O_{\mathbb P^2}(1))$ と書くと，具体的に
\[
\begin{array}{c|cc}
e&c_1(\mathcal A_e)&c_1(\omega_{\mathcal A_e})\\ \hline
2&(-3,-3)&(-5,-5)\\
1&-5H&-7H
\end{array}
\]
なので $\det\omega_{\mathcal A_e}\otimes(\det F)^{-1}=L_e^{-2}$ である．
非零な $F\to\omega_{\mathcal A_e}$ は一般点で非零で，
division rank 1 より一般点で同型となる．その determinant は
$L_e^{-2}$ の非零切断を与えて矛盾する．Serre duality により
$H^2(Z_e,F)=0$ である．一方，$\chi(\mathcal A_e,F)=1$ なので
$H^0(Z_e,F)\ne0$．非零写像 $\mathcal A_e\to F$ の determinant は
非零定数であり，これは全体で同型である．
\end{proof}

\subsection{次数2の exceptional object}

\begin{proposition}
\label{prop:classification2}
$\operatorname{coh}\mathcal A_2$ の局所射影で division rank 1 の
exceptional object は，ちょうど
\[
 T_2^i\mathcal A_2\otimes\mathcal O(j,k)
 \qquad (i,j,k\in\mathbb Z)
\]
である．
\end{proposition}
\begin{proof}
$X_2\subset\mathbb P^2\times\mathbb P^1\times\mathbb P^1$ は
滑らかな ample $(2,1,1)$ hypersurface である．
Weak Lefschetz により $b_2(X_2)=3$ であり，
$H^2$ は因子類で生成される．Hard Lefschetz により $H^4$ も代数的である．
従って Riemann--Roch と Poincar\'e duality から
$\operatorname{rank}K_{\mathrm{num}}(X_2)=1+3+3+1=8$ である．
Kuznetsov の定理 \cite[Theorem~4.2]{KuznetsovQuadric} による
\[
 D^b(X_2)=\langle D^b(Z_2,\mathcal A_2),\,
 p_2^*D^b(Z_2)\otimes\mathcal O_{X_2/Z_2}(1)\rangle
\]
を用いると
$\operatorname{rank}K_{\mathrm{num}}(\mathcal A_2)=8-4=4$ となる．
ここで数値的 Grothendieck 群の分解は，
半直交性と Serre duality により Euler pairing の radical も分解することから従う．

順序付きの対象
$\mathcal A_2,\mathcal P_2,\mathcal A_2(1,0),\mathcal A_2(0,1)$ の
Euler 行列は
\begin{equation}
 G_2=
 \begin{pmatrix}
 1&3&2&2\\1&1&0&0\\0&0&1&-3\\0&0&-3&1
 \end{pmatrix},\qquad\det G_2=16.
 \label{eq:Euler2}
\end{equation}
この計算には
\[
 \chi(\mathcal A_2(a,b))=4ab+a+b+1,\qquad
 \chi(\mathcal P_2(a,b))=4ab+3a+3b+3
\]
と $T_2^{-1}\mathcal A_2=\mathcal P_2(-1,-1)$ を用いる．
従って4つの対象は $K_{\mathrm{num}}(\mathcal A_2)_{\mathbb Q}$ の基底である．

補題\ref{lem:regular-type}より
$c_1(F)=(-3,-3)+2(\alpha,\beta)$，$\alpha,\beta\in\mathbb Z$ と書ける．
rank 条件とこの $c_1$ 条件を基底座標に直すと，ある $z\in\mathbb Q$ について
\[
 [F]=
 \left(1-\frac{\alpha+\beta}{2},\
 z,\ \frac{\alpha-z}{2},\ \frac{\beta-z}{2}\right).
\]
式\eqref{eq:Euler2}は
\[
 \chi(F,F)=1-2\alpha\beta+2z,\qquad
 \chi(\mathcal A_2,F)=1+z+\frac{\alpha+\beta}{2}
\]
を与える．exceptionality から $z=\alpha\beta$，
第2式の整数性から $\alpha\equiv\beta\pmod2$ である．
両者が偶数なら $\mathcal A_2(\alpha/2,\beta/2)$，
奇数なら $\mathcal P_2((\alpha-1)/2,(\beta-1)/2)$ が同じ数値的類を持つ．
この対象を $\mathcal A_2$ に送る自己同値で $F$ を正規化し，
補題\ref{lem:numerical-to-isomorphism}を適用する．
逆に表示中の各対象が局所射影，rank 1，exceptional であることは，
$\mathcal A_2$ の性質と自己同値から従う．
\end{proof}

\subsection{次数1の exceptional object}

まず $\mathcal A_1$ を cubic threefold の Clifford algebra として記述する．
\begin{equation}
 Y=\{ux^2+2vxy+wy^2+2w^2x+2u^2y+2v^3=0\}
 \subset\mathbb P^4_{[x:y:u:v:w]}
 \label{eq:cubic-threefold}
\end{equation}
は直線 $\ell=(u=v=w=0)$ を含む．
ここでは射影束 $\mathbb P(\mathcal V_1)$ を $\mathcal V_1$ 内の
1次元部分空間の射影束とする．射
\[
 \mathbb P_{\mathbb P^2}(\mathcal O^{\oplus2}\oplus\mathcal O(-1))
 \longrightarrow\mathbb P^4,
 \qquad ([u:v:w],[\xi_0:\xi_1:t])
 \longmapsto[\xi_0:\xi_1:tu:tv:tw]
\]
は，$\mathbb P^4\dashrightarrow\mathbb P^2$，
$[x:y:u:v:w]\mapsto[u:v:w]$ のグラフの閉包であり，
$\ell$ に沿う blowup を与える．この射で式\eqref{eq:cubic-threefold}を引き戻し，
例外因子の方程式 $t$ を除くと，狭義変換の方程式は
\[
 u\xi_0^2+2v\xi_0\xi_1+w\xi_1^2
 +2w^2\xi_0t+2u^2\xi_1t+2v^3t^2=0
\]
となる．これは式\eqref{eq:q1}の conic bundle なので
$X_1\simeq\operatorname{Bl}_{\ell}Y$ である．
$X_1$ は滑らかなので $Y\setminus\ell$ も滑らかである．
$\ell$ 上の法線方向の偏微分は $x^2,2xy,y^2$ で，同時には消えない．
従って $Y$ は滑らかな cubic threefold である．

\begin{proposition}
\label{prop:classification1}
$\operatorname{coh}\mathcal A_1$ の局所射影で division rank 1 の
exceptional object は，ちょうど $T_1^n\mathcal A_1$，$n\in\mathbb Z$ である．
\end{proposition}
\begin{proof}
Bernardara--Macr\`i--Mehrotra--Stellari
\cite[Propositions~2.9(ii) and 2.12]{BMMS2012} によれば，
cubic threefold の Clifford algebra について
$[\mathcal B_1],[\mathcal B_0],[\mathcal B_{-1}]$ は
数値的 Grothendieck 群の整数基底である．
本例では $\mathcal B_0=\mathcal A_1$，
$\mathcal B_1=\mathcal P_1$，
$\mathcal B_{-1}=\mathcal P_1(-1)$ である．
$T_1$ を適用して並べ替えると，
$[\mathcal A_1],[\mathcal P_1],[\mathcal A_1(1)]$ が整数基底となる．

すべての $m\in\mathbb Z$ に対して Riemann--Roch より
$p(m):=\chi(\mathcal A_1,T_1^m\mathcal A_1)=1+m(m+1)/2$ である．
従ってこの基底の Euler 行列は
\begin{equation}
 G_1=(p(j-i))_{0\le i,j\le2}
 =\begin{pmatrix}1&2&4\\1&1&2\\2&1&1\end{pmatrix},
 \qquad\det G_1=1.
 \label{eq:Euler1}
\end{equation}
$[F]=(x,y,z)$ と書くと $x,y,z\in\mathbb Z$，
rank 条件は $x+y+z=1$ である．
$n=y+2z$ と置けば
$\chi(F,F)=1+y+4z-n^2$ なので，exceptionality から
\begin{equation}
 (x,y,z)=
 \left(\frac{(n-1)(n-2)}2,\ n(2-n),\
 \frac{n(n-1)}2\right).
 \label{eq:orbit-coordinates}
\end{equation}
この列ベクトルに $G_1$ を掛けると
$(p(n),p(n-1),p(n-2))^{\mathsf t}$ になるため，
式\eqref{eq:orbit-coordinates}は $[T_1^n\mathcal A_1]$ の座標である．
$T_1^{-n}$ で正規化して
補題\ref{lem:numerical-to-isomorphism}を適用する．
逆向きは自己同値から従う．
\end{proof}

\subsection{任意のアーベル圏同値に対する障害}

$B=G/(f)$ を主定理の第3項の algebra とし，
$\pi_B:\operatorname{gr}B\to\operatorname{qgr}B$ を商関手とし，
$\mathcal B=\pi_B B$ と書く．
AS-regular algebra の最小自由分解を双対化した相反公式は
$H_G(t^{-1})=t^\lambda H_G(t)$ を与える．
所定の Hilbert 級数から AS-Gorenstein parameter $\lambda$ は
$5$ または $7$ である．正則正規元による商の change-of-rings 公式から，
$B$ は injective dimension $3$，parameter $\lambda-\deg f=1$ の
AS-Gorenstein algebra となる．

Yekutieli--Zhang \cite[Theorem~4.2 and Corollary~4.3]{YZ1997} を適用すると，
$B$ は $\chi$ を満たし，$\operatorname{qgr}B$ の双対化対象は
$\mathcal B(-1)[2]$ である．同じ局所双対性は
$\mathfrak m=B_{>0}$ に対して
$H_{\mathfrak m}^j(B)=0$ for $j\ne3$ を与え，局所コホモロジーの完全列から
\begin{equation}
 \operatorname{Hom}(\mathcal B,\mathcal B(n))=B_n,\qquad
 \operatorname{Ext}^1(\mathcal B,\mathcal B(n))=0
 \label{eq:saturation}
\end{equation}
を得る．さらに双対性から
$\operatorname{Ext}^2(\mathcal B,\mathcal B)
\simeq\operatorname{Hom}(\mathcal B,\mathcal B(-1))^*=B_{-1}^*=0$
であり，コホモロジー次元が $2$ なので，これより高次の Ext も消滅する．
$\operatorname{End}(\mathcal B)=B_0=k$ と合わせて，
$\mathcal B$ は exceptional である．

仮に $k$-線形アーベル圏同値
$\Phi:\operatorname{qgr}B\xrightarrow{\sim}\mathcal C_e$ が存在するとする．
それは導来圏同値を誘導し，Serre functor と可換である．
$T_B=\Phi^{-1}T_e\Phi$ と置く．
双対性とその表現対象の一意性から
\begin{equation}
 S_B(\mathcal B(n))\simeq\mathcal B(n-1)[2],\qquad
 T_B^n\mathcal B\simeq\mathcal B(n)
 \label{eq:Serre-on-free}
\end{equation}
を得る．ここで第1式は次数移動自己同値が Serre functor と可換であることも用いる．

$F=\Phi(\mathcal B)$ は exceptional である．
閉点に台を持つ任意の単純 $\mathcal A_e$-module $Q$ は
$T_e^2Q\simeq Q$ を満たす．従って $m\gg0$ に対して
\[
 \operatorname{Ext}^1(F,Q)
 \simeq\operatorname{Ext}^1(T_e^{-2m}F,Q)
 \simeq\operatorname{Ext}^1(\mathcal B(-2m),\Phi^{-1}Q)=0.
\]
最後は Yekutieli--Zhang の局所双対性
\cite[Theorem~4.2(1)]{YZ1997} による Serre vanishing である．
実際，$\Phi^{-1}Q=\pi_B M$ と書けば，有限生成次数付き $M$ に対して
\[
 \operatorname{Ext}^1(\mathcal B(-n),\pi_B M)
 =H^2_{\mathfrak m}(M)_n=0\qquad(n\gg0)
\]
となる．最後の消滅は balanced dualizing complex による
局所コホモロジーの右有界性から従う．
$Q$ に値を持つ局所 $\mathcal Hom$ と $\mathcal Ext$ は0次元の台を持つので，
local-to-global spectral sequence によりこれは局所 $\operatorname{Ext}^1$ の消滅でもある．
完備局所 order は semiperfect である．その上の最小射影表示の核が非零なら，
核の非零な単純商があるため，ある単純加群への $\operatorname{Ext}^1$ が非零となる．
従って $F$ は完備化上で射影であり，忠実平坦性により局所射影である．

$F$ の一般点での $\mathscr D$-rank を $r$ とすると，
$\mathcal End_{\mathcal A_e}(F)$ の中心上の rank は $4r^2$ である．
$L_e^2=e$ なので，Riemann--Roch と Serre vanishing は
\[
 \dim\operatorname{Hom}(F,T_e^{2m}F)
 =2er^2m^2+O(m)
\]
を与える．式\eqref{eq:saturation}，\eqref{eq:Serre-on-free}により
左辺は $\dim B_{2m}=2em^2+em+1$ であるから $r=1$ となる．
命題\ref{prop:classification2}，\ref{prop:classification1}により，
$F$ は $T_e$ と可換な自己同値で $\mathcal A_e$ から得られる．

最後に次数2の積を比較する．補題\ref{lem:mu2-surjective}から
$R^{(e)}_2=R^{(e)}_1R^{(e)}_1$ である．
従って $\Gamma(F,T_e)=\bigoplus_{n\ge0}\operatorname{Hom}(F,T_e^nF)$ について
\[
 \dim\frac{\Gamma(F,T_e)_2}
 {\Gamma(F,T_e)_1\Gamma(F,T_e)_1}=0.
\]
一方，式\eqref{eq:Serre-on-free}の対象としての同型を選ぶと，
orbit product の像の次元は合成写像
\[
 \operatorname{Hom}(\mathcal B,\mathcal B(1))
 \otimes\operatorname{Hom}(\mathcal B(1),\mathcal B(2))
 \longrightarrow\operatorname{Hom}(\mathcal B,\mathcal B(2))
\]
の像の次元に等しい．これは式\eqref{eq:saturation}から
$\dim B_1B_1$ である．$G$ の独立な次数2の最小生成元は1個で，
$\deg f>2$ なので
$\dim B_2/(B_1B_1)=1$ となり，矛盾する．
これで主定理の第3項が証明された．


% !TEX root = ../nc-cubic-non-as.tex
\section{Hochschild homology}
\label{sec:hochschild}

次数を $e\in\{1,2\}$ とし，対応する曲面，order，判別曲線を
$Z_e,\mathcal A_e,D_e$ と書く．本節では
\[
 \operatorname{HH}_i(\mathcal C_e)
 :=\operatorname{HH}_i\bigl(\operatorname{Perf}_{\mathrm{dg}}(Z_e,\mathcal A_e)\bigr),
 \qquad \mathcal C_e=\operatorname{coh}\mathcal A_e,
\]
を計算する．$\mathcal A_e$ の局所大域次元は $2$ なので，これは
$D^b(\mathcal C_e)$ の標準的な dg enhancement の Hochschild homology である．
次数の規約は，滑らかな射影多様体 $W$ に対して
\begin{equation}
 \operatorname{HH}_i(W)
 \simeq \bigoplus_{p-q=i} H^q(W,\Omega_W^p)
 \label{eq:hh-hkr-convention}
\end{equation}
とする．

\subsection{判別二重被覆による計算}

$p_e:X_e\to Z_e$ を二次形式 $q_e$ の定める conic bundle とする．
判別曲線 $D_e$ は滑らかであり，その上の退化 conic は相異なる2直線の和である．
従って，退化 conic の一方の既約成分を選ぶ連結な非分岐二重被覆
\[
 \pi_e:\widetilde D_e\longrightarrow D_e,
 \qquad
 (\pi_e)_*\mathcal O_{\widetilde D_e}
 =\mathcal O_{D_e}\oplus\eta_e
\]
がある．ここで $\eta_e$ は第\ref{subsec:residue}節で非自明性を示した2-torsion 線束である．

Kuznetsov の quadric fibration の定理
\cite[Theorem~4.2]{KuznetsovQuadric} を相対次元 $1$ に適用すると，
\begin{equation}
 D^b(X_e)
 =\left\langle
 D^b(Z_e,\mathcal A_e),\,
 p_e^*D^b(Z_e)\otimes\mathcal O_{X_e/Z_e}(1)
 \right\rangle .
 \label{eq:hh-conic-sod}
\end{equation}
$X_e,Z_e,D_e$ が滑らかであることを確認済みなので，
\cite[Theorem~8.12]{KuznetsovHH} の conic bundle に対する公式を適用できる．
\eqref{eq:hh-hkr-convention} の規約に直すと，それは
\begin{equation}
 \operatorname{HH}_i(\mathcal C_e)
 \simeq \operatorname{HH}_i(Z_e)
 \oplus H^{-i}(D_e,\eta_e)
 \oplus H^{1-i}(D_e,\omega_{D_e}\otimes\eta_e)
 \label{eq:hh-discriminant-formula}
\end{equation}
となる．引用文献の homology の添字 $t$ は，ここでは $t=-i$ に対応する．

以下，$g_e=g(D_e)$ と置く．非自明な次数 $0$ の線束 $\eta_e$ は
非零切断を持たず，Serre duality により
\[
 H^0(D_e,\eta_e)=0,
 \qquad
 H^1(D_e,\omega_{D_e}\otimes\eta_e)
 \simeq H^0(D_e,\eta_e^{-1})^*=0.
\]
曲線上の Riemann--Roch から
\begin{equation}
 h^1(D_e,\eta_e)=g_e-1,
 \qquad
 h^0(D_e,\omega_{D_e}\otimes\eta_e)=g_e-1.
 \label{eq:hh-eta-dimensions}
\end{equation}
また $Z_2=\mathbb P^1\times\mathbb P^1$ と $Z_1=\mathbb P^2$ の
Hodge cohomology は対角成分だけなので，
\[
 \operatorname{HH}_i(Z_2)=
 \begin{cases}\mathbb C^4&i=0,\\0&i\ne0,\end{cases}
 \qquad
 \operatorname{HH}_i(Z_1)=
 \begin{cases}\mathbb C^3&i=0,\\0&i\ne0.\end{cases}
\]
次数 $2$ の判別曲線は $(3,3)$ 曲線なので $g_2=4$，
次数 $1$ の判別曲線は平面5次曲線なので $g_1=6$ である．
従って次を得る．

\begin{proposition}
\label{prop:hh-values}
本稿の2つの order に対し，Hochschild homology は
\[
\begin{array}{c|cccc}
 e & \operatorname{HH}_{-1}(\mathcal C_e)
   & \operatorname{HH}_{0}(\mathcal C_e)
   & \operatorname{HH}_{1}(\mathcal C_e)
   & \operatorname{HH}_{i}(\mathcal C_e)\ (|i|\ge2)\\\hline
 2 & \mathbb C^3 & \mathbb C^4 & \mathbb C^3 & 0\\
 1 & \mathbb C^5 & \mathbb C^3 & \mathbb C^5 & 0
\end{array}
\]
である．
\end{proposition}

次数 $1$ と $-1$ の部分は，より具体的には
\begin{equation}
 \operatorname{HH}_1(\mathcal C_e)
 \simeq H^0(D_e,\omega_{D_e}\otimes\eta_e),
 \qquad
 \operatorname{HH}_{-1}(\mathcal C_e)
 \simeq H^1(D_e,\eta_e)
 \label{eq:hh-prym-pieces}
\end{equation}
である．二重被覆の involution を $\iota_e$ とすると，有限射の射影公式から
これらはそれぞれ
\[
 H^0(\widetilde D_e,\omega_{\widetilde D_e})^{-},
 \qquad
 H^1(\widetilde D_e,\mathcal O_{\widetilde D_e})^{-}
\]
に一致する．上付き $-$ は $\iota_e$ の $-1$ 固有空間を表す．
従って非零次数の Hochschild homology は，判別二重被覆の Prym variety
\[
 \operatorname{Prym}(\widetilde D_e/D_e)
 :=\ker\bigl(\operatorname{Nm}_{\pi_e}:
 \operatorname{Pic}^0(\widetilde D_e)\to\operatorname{Pic}^0(D_e)\bigr)^0
\]
の単位元での余接空間と接空間を与えており，その次元は次数 $2$ の例で $3$，
次数 $1$ の例で $5$ である．

\subsection{幾何学的 Euler 数による別計算}

\eqref{eq:hh-conic-sod} と Hochschild homology の加法性
\cite[Corollary~7.5]{KuznetsovHH} を使い，$X_e$ の Hodge 数からも
命題~\ref{prop:hh-values} を確認できる．滑らかな conic の位相的 Euler 数は
$2$，相異なる2直線の和の Euler 数は $3$ なので，
\begin{equation}
 e_{\mathrm{top}}(X_e)
 =2\,e_{\mathrm{top}}(Z_e\setminus D_e)+3\,e_{\mathrm{top}}(D_e)
 =2\,e_{\mathrm{top}}(Z_e)+2-2g_e.
 \label{eq:hh-euler-conic}
\end{equation}
ここでは proper morphism に沿った構成可能関数の Euler 積分を用いている．

次数 $2$ では命題\ref{prop:classification2}の証明で用いた
Weak Lefschetz により $b_1(X_2)=0$，$b_2(X_2)=3$ で，
$H^2$ はすべて $(1,1)$ 型である．随伴公式は
\[
 \omega_{X_2}=\mathcal O_{X_2}(-1,-1,-1)
\]
を与えるので $h^{3,0}(X_2)=0$．Poincar\'e duality と
\eqref{eq:hh-euler-conic} から
\[
 e_{\mathrm{top}}(X_2)=2,\qquad
 b_3(X_2)=(1+3+3+1)-2=6,
 \qquad h^{2,1}(X_2)=h^{1,2}(X_2)=3.
\]
従って $\operatorname{HH}_0(X_2)=\mathbb C^8$，
$\operatorname{HH}_{\pm1}(X_2)=\mathbb C^3$ であり，
$\operatorname{HH}_0(Z_2)=\mathbb C^4$ を除くと命題の第1行となる．

次数 $1$ では $X_1$ は式\eqref{eq:cubic-threefold}の滑らかな cubic threefold $Y$ の直線
$\ell\simeq\mathbb P^1$ に沿う blowup である．
Weak Lefschetz により $b_1(Y)=0$，$b_2(Y)=1$ であり，
$H^2(Y)$ は超平面類で生成される．blowup の cohomology 公式は
$H^2$ と $H^4$ に各1次元を加え，$H^3$ には
$H^1(\mathbb P^1)=0$ を加える．従って
$b_1(X_1)=b_5(X_1)=0$，$b_2(X_1)=b_4(X_1)=2$ であり，
$H^2(X_1)$ と $H^4(X_1)$ はそれぞれ $(1,1)$ 型と $(2,2)$ 型である．
また $\omega_Y=\mathcal O_Y(-2)$ と blowup の公式から
$h^{3,0}(X_1)=0$ である．よって
\[
 e_{\mathrm{top}}(X_1)=-4,\qquad
 b_3(X_1)=(1+2+2+1)-(-4)=10,
 \qquad h^{2,1}(X_1)=h^{1,2}(X_1)=5.
\]
従って $\operatorname{HH}_0(X_1)=\mathbb C^6$，
$\operatorname{HH}_{\pm1}(X_1)=\mathbb C^5$ であり，
$\operatorname{HH}_0(Z_1)=\mathbb C^3$ を除くと命題の第2行となる．

\subsection{帰結}

\begin{corollary}
いずれの $D^b(\mathcal C_e)$ も full exceptional collection を持たない．
また，次数 $e$ の滑らかな可換 del Pezzo 曲面の導来圏とは同値でない．
\end{corollary}

\begin{proof}
exceptional object の生成する圏の Hochschild homology は次数 $0$ の
$\mathbb C$ だけなので，加法性により full exceptional collection を持つ圏では
$\operatorname{HH}_{i}=0$ for $i\ne0$ となる．これは
命題~\ref{prop:hh-values} の $\operatorname{HH}_{\pm1}\ne0$ と矛盾する．
導来同値の非存在は数値的 Grothendieck 群からも分かる．
第\ref{sec:unpointed}節で求めたその rank は，次数2の例で $4$，次数1の例で $3$ である．
一方，滑らかな可換 del Pezzo 曲面 $S$ では Kodaira vanishing より
$H^i(S,\mathcal O_S)=0$ for $i>0$ で，Noether の公式は
$e_{\mathrm{top}}(S)=12-K_S^2=12-e$ を与える．
Hodge cohomology は対角成分だけで，$H^2$ は因子類で生成されるから，
$\operatorname{rank}K_{\mathrm{num}}(S)=12-e$，すなわち $10$ または $11$ である．
Euler pairing と数値的 Grothendieck 群は三角圏同値で保たれるため，導来同値は存在しない．
\end{proof}




\section{成果の比較と残る課題}
量子Fermat点と符号型族の奇数次Hochschild homologyはOU型AS surfaceへの障害であるが、AS四次元ambientの超曲面表示自体を否定しない。
次数三Clifford例では$\dim Q(R)_3=19$と標準AS四次元代数の$20$の差が、構造対象を忘れた非表示の核心である。
次数一・二Clifford例は最小生成次数を固定した非表示であり、別のambientによる肯定的な表示も同時に得る。
現在の課題は、次数三Clifford例について非周期的$\mathbb Z$-algebraまで障害を拡張できる条件の確定、符号型の相対変形とOU型部分の位相、及びこれらの明示的結果の初出の徹底調査である。
今回の統合によって新たな証明が自動的に得られたとは主張しない。

\appendix
% !TEX root = ../nc-cubic-non-as.tex
\section{厳密計算の入力・方法・出力}
\label{sec:computation}

本付録の計算は Python 3 と SymPy を用いる有理数体上の厳密計算であり，
浮動小数点近似を用いない．

\subsection{判別曲線の滑らかさ}

次数2の入力は式\eqref{eq:q2}の determinant $\det q_2$ である．
$s_i=t_j=1$，$(i,j)\in\{0,1\}^2$ とする4チャートで，
\[
 (s_0,s_1)=
 \begin{cases}(1,x)&i=0,\\(x,1)&i=1,\end{cases}
 \qquad
 (t_0,t_1)=
 \begin{cases}(1,y)&j=0,\\(y,1)&j=1\end{cases}
\]
を同時代入して $F_{ij}\in\mathbb Q[x,y]$ を得る．
次数1の入力は式\eqref{eq:q1}の determinant $\det q_1$ であり，
$u=1,v=1,w=1$ の各チャートで残る2変数の多項式 $F$ を得る．
それぞれ
\[
 I=(F,\partial_xF,\partial_yF)
\]
の辞書式順序の Gröbner 基底を計算すると，出力はすべて $\{1\}$ である．
従って各 $I$ は単位イデアルであり，$\mathbb C$ 上にも共通零点を持たない．
7チャートが曲面を被覆するので，これは判別曲線全体の滑らかさを証明する．

\subsection{Hilbert 級数と Euler pairing}

スクリプトは偶奇の Hilbert 関数を $m$ の多項式として展開し，
本文の式との差が零多項式であることを確認する．
Hilbert 級数の比較も有理関数の差を通分・約分し，
恒等的に零になることを確認する．従って有限個の係数の一致からの推測ではない．
また式\eqref{eq:Euler2}，\eqref{eq:Euler1}の determinant と，
rank 1 の自己 Euler 数，式\eqref{eq:orbit-coordinates}の
全整数 $n$ に関する多項式恒等式を確認する．

\subsection{Hochschild homology の数値的検証}

第\ref{sec:hochschild}節の証明に現れる種数，Euler 数，
Betti 数，Hochschild homology の次元を整数として確認する．
別に Chern class の積分によって Euler 数を計算する．
次数2では
\[
 A^\bullet(\mathbb P^2\times\mathbb P^1\times\mathbb P^1)_{\mathbb Q}
 =\mathbb Q[h,a,b]/(h^3,a^2,b^2),\qquad [X_2]=2h+a+b
\]
において
\[
 c(TX_2)=\frac{(1+h)^3(1+a)^2(1+b)^2}{1+2h+a+b}\bigg|_{X_2},
 \qquad
 c_3(TX_2)=2abh-ah^2-bh^2.
\]
$[X_2]c_3(TX_2)$ の $h^2ab$ の係数は $2$ である．
次数1の cubic threefold については
\[
 c(TY)=\frac{(1+h)^5}{1+3h}\bigg|_Y,\qquad
 c_3(TY)=-2h^3,\qquad \int_Yh^3=3,
\]
従って $e_{\mathrm{top}}(Y)=-6$ となる．直線に沿う blowup で
$e_{\mathrm{top}}(\mathbb P^1)=2$ が加わり，
$e_{\mathrm{top}}(X_1)=-4$ を得る．
これは conic bundle のファイバーによる Euler 数の計算と一致する．

\Needspace{75mm}
\subsection{再現方法}

repository の root から
\begin{verbatim}
python3 -m pip install -r requirements.txt
python3 scripts/verify_all.py
python3 scripts/build.py --papers
\end{verbatim}
を実行する．検証対象は次の3ファイルである．
\begin{verbatim}
scripts/del_pezzo_degree2_verification.py
scripts/nc_dp1_quintic_checks.py
scripts/hochschild_verification.py
\end{verbatim}
この三つに加え，検証の共通入口は同梱の次数三の二つのスクリプトも実行する．
成功時の出力はその場で再生成し，重複する実行ログは収録しない．
PDF の組版には LuaLaTeX，BibTeX と日本語フォントを用いる．
必要な環境と Cursor からの実行方法は \texttt{README.md} に記す．


\bibliographystyle{amsplain}
\bibliography{references}
\end{document}
